% Modernised reading — fonds Alexandre Grothendieck, shelfmark 134-2, pages
% 1-80, which is all of the folder that is transcribed: batches 1 to 4 of a
% folder of 184 pages (issue #44). Pages 81-184 are not transcribed and are not
% read here.
%
% This is an INTERPRETATION, not a transcription. It re-reads the pages in
% today's notation and vocabulary, states the hypotheses the manuscript leaves
% implicit, and drops the critical apparatus so that the mathematics reads
% straight through. Everything the brackets and underlines carried moves into
% a footnote. In any dispute of substance the transcription governs: whoever
% cites Grothendieck cites batch-NN.fr.tex.
%
% Scope, and what is deliberately NOT restated. Under the special rules of
% this folder (#44), the letter to Daniel Quillen of 19-23.2.1983 and the
% numbered sections of Pursuing Stacks (pages 14-36), and the typescript
% « NOTES to Chapter I and Appendix » (pages 78-80), are only noted in the
% transcription, because Maltsiniotis's edition (À la poursuite des champs,
% vol. I, SMF, Documents mathématiques 20, 2022) prints them. This reading
% cites that edition for them and does not paraphrase them. What it reads is:
% the plans and tables of contents (pp. 1-12), the dedication (p. 13), his
% 1983 introduction to the Appendix (p. 37), and his three letters to Lawrence
% Breen of 5.2.1975 (pp. 38-47), 17.2.1975 (pp. 48-55) and 17/19.7.1975
% (pp. 56-77), which the printed volume I does not include. Breen's own
% letters are not in these pages and are not reconstructed: only what
% Grothendieck's letters say of them is reported.
%
% Notation fixed for the whole document: S = infinity-groupoids (spaces);
% S_{<=n} = n-groupoids = homotopy n-types; an « n-champ » on a topos X is an
% n-truncated object of the infinity-topos of sheaves of spaces on X;
% Pi_n(X) = tau_{<=n} of the shape (pro-object in general); his cursive
% « espace classifiant » functor (read T in the transcription) is written B
% here, and his « n-topos classifiant » T_n keeps its letter. A map f is
% « k-connexe » if it is bijective on pi_i for i < k and surjective on pi_k.
%
% Substantive departures from the pages, all footnoted where they occur:
%  - p. 48: the chain « critère d'Artin-Mazur <=> critère des n-gerbes <=> (A)
%    <=> (B) » is off by one degree; the three conditions are stated
%    separately, with their exact homotopical meaning and counterexamples.
%  - p. 54: (B') as written is false in the direction « f_n n-équivalence =>
%    ... »; corrected to (n+1)-connectedness of f_{n+1}.
%  - p. 50-51: the expectation that van Kampen fails for Pi_2 is not borne out
%    (truncation is a left adjoint); stated with the hypotheses that make it
%    true.
%  - p. 53: B^2 F needs a braided (E_2) structure, and all B^k need a Picard
%    (E_infinity) structure, not a strict one.
%  - p. 52: the « morphisme canonique de (n-1)-groupoïdes » is between
%    (n+1)-groupoids.
%  - p. 45: Lf_* versus Rf_* — the direct image of simplicial sheaves is
%    derived on the right.
%  - p. 66: Gamma_k on the left of the arithmetic local duality is Gamma_K;
%    D_S(i^* F) is D_s(i^* F).
%  - p. 67: « the construction in (B) of the local jacobian complex » is (C).
%  - p. 68: target (n-Cat) made (n-1)-Cat, consistent with (n-1)-systems.
% The modern situating (n-stacks, homotopy hypothesis, shape, exodromy,
% six-functor formalisms, flat duality) is ours; every such passage says so.
%
% Pass: Opus 5.5 (claude-opus-5-5), 2026-10-03 — first pass, unchecked against
% the transcription by a human.
\documentclass[11pt,a4paper]{article}
\input{../preamble/grothendieck}

\folder{134-2}
\pages{1}{80}
\dating{1975-[1983]}
\watermark{Édition de démonstration\\interprétation personnelle de l'œuvre}
\foldertitle{[Chapitre I : Take off (pages 1 à 65) et table des matières provisoire] : tapuscrits annotés (19/02-22/02/1983), note manuscrite (s.d.), copies de lettre (1975, s.d.) — lecture modernisée du dossier entier}

\begin{document}

\section*{Résumé}

\begin{resume}
Ce dossier est l'ouverture d'\emph{À la poursuite des champs}
(\emph{Pursuing Stacks}), le long texte que Grothendieck commence en février
1983 sous la forme d'une lettre à Daniel Quillen. Il en contient trois
choses : les plans et tables des matières de l'ouvrage tel qu'il le voyait
alors ; le premier chapitre, « Take-off », qui est la lettre à Quillen ; et
l'appendice de ce chapitre, trois lettres de 1975 à Lawrence Breen que
Grothendieck a fait retaper pour les joindre au volume.

Il faut dire d'emblée ce que cette lecture couvre. Le dossier compte 184
pages ; seules les pages 1 à 80 sont transcrites, et c'est d'elles seules
qu'il est question ici. Parmi elles, la lettre à Quillen et les sections
numérotées (pages 14 à 36) et les notes au chapitre (pages 78 à 80) ne sont
pas reprises : elles sont imprimées dans l'édition de Georges Maltsiniotis
(\emph{À la poursuite des champs}, vol.~I, SMF, 2022), à laquelle on renvoie.
Ce qui est lu, ce sont les tables (pages 1 à 12), une dédicace (page 13), et
surtout les trois lettres à Breen, du 5 février, du 17 février et des 17--19
juillet 1975 (pages 37 à 77), que ce volume imprimé ne contient pas.

De quoi ces lettres parlent-elles ? D'une intuition qui vient d'arriver :
qu'un « espace à déformation près » et une certaine structure purement
algébrique --- un $n$-\emph{groupoïde} --- sont au fond la même chose. Un
groupoïde, c'est une collection d'objets et de flèches toutes inversibles.
Si l'on prend pour objets les points d'un espace et pour flèches les chemins
à déformation près, on obtient le groupoïde fondamental, qui retient les
composantes connexes et le groupe fondamental, et oublie le reste. Pour
retenir davantage, il faut cesser de passer au quotient trop tôt : garder les
chemins eux-mêmes, puis les déformations entre chemins, puis les déformations
entre déformations, sur $n$ étages. C'est un $n$-groupoïde, et l'idée de 1975
est qu'il retient exactement la forme de l'espace jusqu'en dimension $n$.

Quel problème cela résout-il ? Grothendieck répond à Breen, qui lui demandait
des exemples « convaincants » d'objets de ce genre, et lui en donne : les
gerbes, les algèbres d'Azumaya, les champs de Picard. Mais l'argument de fond
est ailleurs. L'algèbre homologique classique travaille avec des complexes
de groupes \emph{abéliens}, et beaucoup de théorèmes de géométrie (ceux de
Lefschetz, par exemple) se démontrent aussi pour des coefficients qui ne
commutent pas, où l'outil manque. Ce qui joue pour les objets non
commutatifs le rôle des complexes, ce seraient ces $n$-groupoïdes, et leurs
versions qui varient d'un point à l'autre d'un espace : les
$n$-\emph{champs}. La cohomologie d'un espace devient alors une opération
très simple --- prendre les sections globales d'un champ --- sans résoudre ni
l'espace ni les coefficients.

La troisième lettre est plus vaste. Elle esquisse un programme de dualité
conçu à la fin des années cinquante, où la théorie du corps de classes,
locale et globale, devient une formule de dualité entre deux cohomologies ;
et elle esquisse une « topologie modérée » dans laquelle un espace muni d'une
stratification se décrit par un diagramme fini de types d'homotopie, d'où
un « type d'homotopie fin » d'un espace ou d'un schéma.

Une analogie aide à tenir ensemble ces fils. Un revêtement d'un espace, ce
sont des données qui sont constantes là où l'on regarde de près, mais qui
peuvent se tordre quand on fait le tour d'un trou. Grothendieck pose que
l'espace, vu « à homotopie près », n'est rien d'autre que la collection de
toutes ces données localement constantes, à tous les étages ; et que si l'on
autorise ces données à changer de nature le long d'une stratification, on
voit apparaître un type d'homotopie plus fin, qui retient comment les
strates sont recollées.

Où cela mène-t-il ? Presque tout ce que ces lettres demandent a reçu depuis
un nom, parfois une démonstration. L'hypothèse que les $n$-groupoïdes
modélisent les $n$-types d'homotopie est l'« hypothèse d'homotopie » ; les
$n$-champs sont les objets tronqués d'un $\infty$-topos, et les
« $n$-topos » qu'il redoutait existent ; le type d'homotopie d'un topos est
sa « forme » (\emph{shape}) ; la description d'un espace stratifié par
diagramme est la théorie des chemins de sortie, aujourd'hui
« exodromie ». Les situer est le travail de cette lecture, qui signale
chaque fois ce qui est son propre commentaire. Quelques-unes de ses
affirmations de 1975 sont décalées d'un degré, et on le dit à l'endroit où
elles se trouvent.

Reste la manière, qui se lit sans aucun bagage. À Breen qui lui objecte que
la classification des objets géométriques se ramène finalement à des
invariants cohomologiques bien connus, il répond que c'est négliger la
différence entre \emph{comprendre} un objet et connaître sa classe à
isomorphisme près. Qu'on ait jusqu'ici toujours pu éluder ces animaux
n'est pas pour lui une raison de continuer : beaucoup de mathématiques
« sérieuses » n'existeraient pas si l'on n'avait éclairci que ce qu'on était
\emph{forcé} d'éclaircir. Et devant une définition technique qui marche mais
n'éclaire rien, il propose une autre méthode : dresser, comme une liste de
cadeaux de Noël, la liste complète de ce qu'on voudrait avoir, et attendre
qu'elle désigne d'elle-même la construction --- « you just follow your nose ».

\keywords{n-stack, n-groupoid, homotopy hypothesis, fundamental n-groupoid, Picard stack, gerbe, Brauer group, crossed module, shape of a topos, infinity-topos, six-functor formalism, flat duality, geometric class field theory, exodromy, stratified spaces, test category}
\end{resume}

\section*{Le fil du dossier, et ce que cette lecture couvre}

\subsection*{Les stations}

Le dossier est un tapuscrit de 1983 annoté de sa main, dans l'ordre où
Grothendieck l'a monté pour en faire le premier chapitre de \emph{Pursuing
Stacks}. On y distingue sept stations.

\begin{itemize}
  \item[1.] \emph{Pages 1 à 12} : la couverture, « The Modelizing Story », le
  plan en parties I à X et deux états successifs de la table des matières,
  §§~1 à 140, plus la table de l'appendice (App.~1 à 18). Ce sont ses tables
  de travail ; elles sont lues ici comme un document.
  \item[2.] \emph{Page 13} : une dédicace manuscrite.
  \item[3.] \emph{Pages 14 à 36} : le chapitre I, « Take off », c'est-à-dire
  la lettre à Daniel Quillen des 19--23 février 1983, §§~1 à 13. Imprimé par
  Maltsiniotis ; non repris ici.
  \item[4.] \emph{Page 37} : l'avertissement de 1983 à l'appendice, en
  anglais.
  \item[5.] \emph{Pages 38 à 47} : première lettre à Breen, Villecun,
  5.2.1975 (App.~1 à 5).
  \item[6.] \emph{Pages 48 à 55} : deuxième lettre, Villecun, 17.2.1975
  (App.~6 à 8).
  \item[7.] \emph{Pages 56 à 77} : troisième lettre, Villecun, 17/19 juillet
  1975, dans une traduction anglaise revue par lui (App.~9 à 18).
\end{itemize}

Suivent, pages 78 à 80, les premières pages des « NOTES to Chapter I and
Appendix », imprimées dans la même édition et non reprises ici.

\subsection*{L'épine dorsale}

Le dossier ne construit pas un objet de page en page ; il en \emph{annonce}
un. Les tables disent où \emph{Pursuing Stacks} devait aller, la lettre à
Quillen dit par où il commence, et les lettres à Breen disent d'où il vient :
de l'intuition de 1975 que les $n$-groupoïdes non stricts modélisent les types
d'homotopie $n$-tronqués --- l'avertissement de la page 37 le dit en propres
termes. Les trois lettres déploient cette intuition dans un ordre qui est
celui de cette lecture :
\begin{enumerate}
  \item des exemples de $2$-catégories de Picard (pages 38 à 40) ;
  \item les $n$-champs comme cadre naturel d'une algèbre homologique non
  commutative (pages 40 à 42) ;
  \item le dictionnaire entre $n$-groupoïdes et types d'homotopie tronqués,
  absolu puis relatif à un topos (pages 42 à 47) ;
  \item le $n$-groupoïde fondamental d'un topos et la classification de ses
  systèmes locaux (pages 48 à 51) ;
  \item la cohomologie comme intégration de champs (pages 51 à 55) ;
  \item en juillet, le prolongement vers la dualité, la topologie modérée et
  la linéarisation des types d'homotopie (pages 56 à 77).
\end{enumerate}

\subsection*{Conventions, valables pour tout le dossier}

\begin{itemize}
  \item On note $\mathcal{S}$ l'$\infty$-catégorie des espaces, ou
  $\infty$-groupoïdes, et $\mathcal{S}_{\leqslant n}$ celle des
  $n$-groupoïdes, c'est-à-dire des $n$-types d'homotopie (espaces dont les
  $\pi_i$ sont nuls pour $i > n$ en tout point base). Ce dictionnaire est
  \emph{l'hypothèse d'homotopie} ; on dit à la page 43 sous quelle forme elle
  est un théorème.
  \item Un \emph{$n$-champ} (en groupoïdes) sur un topos $X$ est un objet
  $n$-tronqué de l'$\infty$-topos $\mathrm{Sh}(X;\mathcal{S})$ des faisceaux
  d'espaces sur $X$ ; pour $n = 0$ ce sont les faisceaux, pour $n = 1$ les
  champs en groupoïdes de Giraud.\footnote{Pour des objets tronqués, la
  différence entre cet $\infty$-topos et son hypercomplété ne joue pas ;
  elle reviendra à la page 44, à propos des $\infty$-catégories.}
  Un \emph{$n$-champ de Picard strict} est l'objet associé à un complexe de
  faisceaux abéliens concentré sur $n+1$ degrés consécutifs ; un
  \emph{champ de Picard} au sens de Deligne (SGA~4, XVIII, 1.4) est le cas
  $n=1$.
  \item $\Pi_n(X)$ désigne le $n$-groupoïde fondamental de $X$ --- la
  troncature en degré $n$ de sa \emph{forme}, en général un pro-objet. La
  lettre du 17 février le note d'abord $\mathcal{C}_n$ ; les deux notations
  désignent le même objet.
  \item Le foncteur « espace classifiant » des pages 52 et 53, qu'il note
  d'une cursive lue $\mathcal{T}$ dans la transcription, est noté ici $B$,
  comme aujourd'hui ; la lettre $\mathcal{T}_n$ reste réservée au
  « $n$-topos classifiant » de la page 50.
  \item Une application $f$ est dite \emph{$k$-connexe} si elle induit une
  bijection sur les $\pi_i$ pour $i < k$ et une surjection sur $\pi_k$, en
  tout point base. C'est la convention dont dépendent les corrections des
  pages 48 et 54.
  \item Les intertitres « App.~1 » à « App.~18 » sont de sa main, ajoutés en
  1983 ; ils servent ici à repérer les passages.
\end{itemize}

\subsection*{Ce que le dossier n'établit pas, et ce qui manque}

\begin{itemize}
  \item La lettre de juillet passe de sa page 1 (page 56 du dossier) à sa
  page 12 (page 57) : ses pages 2 à 11 ne sont pas dans le dossier
  transcrit.\footnote{Sa pagination de 1983 va de 40 à 42 sur ces deux
  feuillets ; elle suggère qu'un seul feuillet de son tapuscrit manque, la
  pagination propre de la lettre en suggère dix. Le passage manquant
  contenait sans doute ses réponses aux questions 1 à 6 de Breen --- la
  page 67 répond à la « question 7 » ---, mais ce n'est qu'une inférence.}
  \item Le complexe $X(M)$ de la page 39 n'est pas défini sur ces pages.
  \item Les commentaires de 1983 annoncés page 37, « notes (22) à (31) »,
  sont dans les « NOTES » des pages 78 sq., imprimées par Maltsiniotis et
  non reprises ; leur suite est au-delà de la page 80.
  \item Les lettres de Breen ne sont pas dans ces pages, et on ne les
  reconstruit pas.
  \item Les pages 81 à 184 ne sont pas transcrites ; on ne sait de leur
  contenu que ce qu'en dit l'inventaire.
\end{itemize}

\pagerange{1}{12}
\section*{Le plan de \emph{Pursuing Stacks} (pages 1 à 12)}

La couverture porte au crayon, et seulement cela, « The Modelizing Story ».
La page suivante donne un plan en parties, tapé puis remanié à la main\footnote{Page 2. Les numéros sont corrigés à la main (« VII » devient
« VIII », « VIII » devient « IX », « X » est récrit sur un « VIII » biffé) ;
la partie VII est ajoutée verticalement dans la marge gauche ; « and test
functors », « and contractibility structures », « and canonical modelizers »
sont biffés dans les intitulés de II, III et IV, et l'intitulé primitif de V
était « Homology and cohomology (abelianization of homotopy types) ». Un mot
au crayon, à peine lisible, est lu avec doute « homotopy types ».} :
\begin{itemize}
  \item[I] Take-off ;
  \item[II] Test categories ;
  \item[III] Homotopy structures ;
  \item[IV] Asphericity structures ;
  \item[V] Abelianization ;
  \item[VI] Schematization ;
  \item[VII] Abelianization (II) ;
  \item[VIII] Homotopy properties of (Cat) and $A^{\wedge}$ (closed model
  structures) ;
  \item[IX] Derivators ;
  \item[X] Back to topoi.
\end{itemize}
Sous la liste, un schéma de dépendances relie les parties ; on le reproduit
tel qu'il est tracé, la partie I restant isolée\footnote{Deux flèches de ce schéma sont biffées : l'une de V vers un « VI »
biffé placé à sa droite, l'autre de V vers VIII. La partie X n'y figure pas.} :

\begin{tikzcd}[column sep=small, row sep=small]
\mathrm{I} & \mathrm{II} \arrow[r, Rightarrow] & \mathrm{III} \arrow[r, Rightarrow] & \mathrm{IV} \arrow[dd, Rightarrow] & \mathrm{V} \arrow[d, Rightarrow] \arrow[dr] & \\
 & & & & \mathrm{VII} & \mathrm{VI} \\
 & & & \mathrm{VIII} \arrow[d, Rightarrow] & & \\
 & & & \mathrm{IX} & &
\end{tikzcd}

Les pages 3 à 12 donnent deux états de la table. Le premier (page 3, sa
page « i ») titre encore « The Modilizing Story » et attribue le chapitre I à
« A letter to … », sans nom ; le second (page 4) porte le titre général
« PURSUING STACKS (À la poursuite des champs). First episode: THE MODELIZING
STORY (Histoire de Modèles) » et précise « a letter to Daniel Quillen ». La
table de l'appendice --- les dix-huit intertitres des lettres à Breen ---
n'est qu'au premier état ; elle concorde avec les marques « App. » portées
dans les marges des lettres, qui sont donc de la même campagne de 1983.\footnote{Le premier état comporte une consigne d'impression : « Les chiffres
soulignés 5, 9, 13, App.~3, App.~6, App.~11, App.~12 etc sont à imprimer en
caractères gras ». Au second état, ce sont des numéros cerclés à l'encre
(5, 9, 13, puis 23, 28, 31, 33, 37, 46, 50, 51, 62, 72, 75, 77, 79, 81, 102,
109 à 111, 118, 119, 124, 131, 132) : ce sont les sections qu'il tenait pour
principales. Le titre de l'App.~6, « Champs essentiellement localement
constants et types d'homotopie », déborde le bord de la feuille numérisée et
reste incomplet.}
La transcription note que ces tables de travail diffèrent de la table
imprimée ; on ne les a pas collationnées avec elle.

Ce qu'elles montrent, c'est la forme du projet à la date des pages 7 et 8 :
§§~1 à 66 en deux parties (« Take-off », puis « Test categories and test
functors » et « Homotopy structures, contractibility structures »), une
partie IV refaite au propre à la main après qu'une première version tapée a
été barrée (§§~67 à 99), puis abélianisation (§§~100 à 109),
schématisation (§§~110 à 133) et linéarisation (§§~133 à 140), la table
s'arrêtant sur « (to be completed) ».\footnote{Page 7, en marge du § 65 :
« p. 168 -- » suivi d'un mot lu avec doute « 199 », renvoi à la pagination
du tapuscrit. Page 10, la place de l'insertion « pseudo-topoi » au § 104 est
incertaine. Le § 133, « Birth of Suleyman », figure deux fois, en fin de VI
et en tête de VII.}

Quelques noms d'aujourd'hui pour s'orienter dans ces intitulés --- ils sont
de nous, et l'ouvrage, qui est de 1983, ne pouvait pas les connaître, sauf
le premier :
\begin{itemize}
  \item les « catégories test » et « modélisateurs » (§§~26 à 44, 65) et le
  « basic localizer » $\underline{W}$ (§§~60 à 65) sont devenus, chez
  Maltsiniotis (\emph{La théorie de l'homotopie de Grothendieck}, Astérisque
  301, 2005) et Cisinski (\emph{Les préfaisceaux comme modèles des types
  d'homotopie}, Astérisque 308, 2006), la théorie des catégories test et des
  localisateurs fondamentaux ; Cisinski y démontre la conjecture de
  Grothendieck que le localisateur fondamental des équivalences faibles
  usuelles de $(\mathrm{Cat})$ est le plus petit ;
  \item le § 87 commente l'article de Thomason (1980) qui munit
  $(\mathrm{Cat})$ d'une structure de catégorie de modèles fermée --- c'est
  le seul de ces noms que le texte cite lui-même ;
  \item les « dérivateurs » (§ 69, partie IX) sont la théorie qu'il
  développera ensuite (\emph{Les Dérivateurs}), parallèlement à Heller ;
  \item la « schématisation » des types d'homotopie (§§~110 à 133) est ce
  que Toën a construit sous le nom de champs affines (2006) et de types
  d'homotopie schématiques (Katzarkov, Pantev, Toën, 2008) ;
  \item la paire $\widetilde{LH}_\bullet \dashv \widetilde{K}$ du § 121
  entre types d'homotopie abéliens et types d'homotopie pointés est, sauf
  erreur de lecture du titre, l'adjonction entre homologie réduite (le
  groupe abélien libre engendré, à la Dold--Thom) et l'oubli, via
  Dold--Kan.
\end{itemize}

\pagerange{13}{13}
\section*{La dédicace (page 13)}

La page est manuscrite. Une première dédicace, « À ceux qui furent mes
élèves --- à qui j'ai donné du meilleur de moi-même, et aussi du pire … »,
est barrée de traits obliques. Celle qui reste est à son frère, Dodek
Shapiro, qu'il aurait « voulu rencontrer un jour », et dont les seuls signes
tangibles sont « quelques lettres et photos d'enfant, anciennes de près d'un
demi-siècle », datées d'Oulianovsk et adressées à leur père,
Sascha.\footnote{Le texte dit « datées d'Ulianovsk en Sibérie » ; Oulianovsk
est sur la Volga, en Russie d'Europe, et non en Sibérie. On rapporte la page
sans la corriger. Un mot biffé et noirci sous « Sibérie » est lu avec doute
« maman » ; « notre » père est corrigé en « son » père, « (qui est aussi le
mien) ».}

\pagerange{14}{36}
\section*{Le chapitre I : la lettre à Quillen (pages 14 à 36)}

Les pages 14 à 36 sont le tapuscrit annoté de la lettre à Daniel Quillen
commencée aux Aumettes le 19 février 1983 et poursuivie jusqu'au 23 (avec
une note marginale du 25), qui forme les §§~1 à 13 de \emph{Pursuing Stacks}
et se termine par un post-scriptum ; un bandeau collé page 14 l'intitule
« Chapter I : TAKE OFF / A letter to …. / (followed by an Appendix: Three
letters to Larry Breen) ». Ce texte est imprimé dans G.~Maltsiniotis (éd.),
\emph{À la poursuite des champs}, vol.~I, SMF, Documents mathématiques 20,
2022, et c'est là qu'il faut le lire ; on ne le restitue ni ne le paraphrase
ici. Les titres de ses sections sont ceux des tables des pages 3 et 4 : de
« The importance of innocence » (§ 1) à « An urgent reflection on proper
names: “stacks” and “coherators” » (§ 13).\footnote{La transcription relève
seulement, pour s'y repérer, la correspondance des pages : pages 14 à 20,
§§~1 à 4 (sa p.~3 classée après sa p.~4) ; pages 21 à 36, §§~4 (fin) à 13,
avec les dates « 21.2. » (page 28), « 22.2.1983 » (page 34) et « PS (23.2) »
(page 36).}

\pagerange{37}{37}
\section*{L'avertissement de 1983 à l'appendice (page 37)}

L'appendice s'ouvre sur une page en anglais, qui dit d'où viennent les
textes qui suivent et ce qu'il en a fait. Les deux premières lettres sont
reproduites telles quelles, en français ; pour la deuxième il se sert d'une
copie dactylographiée faite par Breen en 1975. La troisième paraît dans une
traduction anglaise faite l'année précédente par Ronnie Brown, avec l'aide
de Breen et de J.-L.~Loday, à partir d'une copie à peine lisible de la
lettre manuscrite ; il dit s'en servir plutôt que de l'original, qu'« aucun
imprimeur ne pourrait déchiffrer », et l'avoir corrigée là où son écriture
avait égaré les traducteurs. Il a sauté le début de la première lettre, ajouté
les intertitres App.~1 à 18 et des commentaires, les notes (22) à
(31).\footnote{Il en résulte que les mots de la troisième lettre, telle
qu'on la lit pages 56 à 77, ne sont pas les siens mais ceux d'une
traduction qu'il a revue. On la cite donc en anglais, entre guillemets,
quand une formule compte, sans la retraduire.}

Il situe ensuite les lettres. Les deux premières tentaient d'expliquer à
Breen « some of the main points of the programme I had in mind around the
notions of $n$-categories and $n$-stacks », sous l'impulsion d'une intuition
neuve pour lui : « that (non strict) $n$-groupoids should model (in a
suitable sense) $n$-truncated homotopy types ». La troisième, réponse aux
questions de Breen, est « of a wider scope » : un programme de dualité
remontant à la fin des années cinquante, y compris une formulation
cohomologique de la théorie du corps de classes « géométrique » locale et
globale, qui « appears here for the first time in print » ; puis le besoin
d'une « tame topology » pour une théorie du dévissage des espaces
stratifiés et des voisinages tubulaires étales, en vue d'un « fine homotopy
type » défini par le système projectif des types d'homotopie indexés
associés aux stratifications équisingulières.

\pagerange{38}{47}
\section*{Première lettre à Breen, Villecun, 5 février 1975 (pages 38 à 47)}

La lettre est datée « Villecun 5.2.1975 » et commence, dans le volume, par
« Cher Breen » suivi d'une coupure. Le début sauté transparaît sous la bande
collée : il y remercie Breen d'une longue lettre et parle de sa propre
lettre à Deligne et d'une interprétation, due à Breen, d'un groupe
d'obstruction comme un $\mathrm{Ext}$.\footnote{Page 38. Ce début est barré
de grandes croix et ne se lit qu'à travers le papier ; la transcription en
donne des fragments, dont plusieurs mots incertains (« Ton scepticisme »,
« était fondé », « cohomologique », « Ext (X(M),T) ») et des lacunes. On n'en
tire rien de plus que ce qui est dit ici. La « lettre à Deligne » est très
vraisemblablement celle du 27 août 1974 que contient le dossier 134-1, où
paraît le « drôle de $\mathrm{Ext}^2$ » que mentionne la page 39.}
Que les pages 41 à 47 soient la fin de cette même lettre résulte de la
continuité : la phrase interrompue au bas de la page 40 (« des théorèmes de
Lefschetz ») reprend en haut de la page 41, et les intertitres App.~2 à 5 se
suivent.\footnote{La transcription de ces pages, faite lot par lot, laissait
le destinataire ouvert pour les pages 41 à 47 ; lu d'un seul tenant, le
dossier le fixe, et l'avertissement de la page 37 compte bien deux lettres en
français avant celle de juillet.}

\pagerange{38}{40}
\subsection*{App.~1 : des $2$-catégories de Picard (pages 38 à 40)}

Breen espérait, dit-il, qu'on lui prouverait que les $2$-catégories de
Picard, les $n$-catégories « et autre faune de ce genre » sont
indispensables à des mathématiques sérieuses. Grothendieck craint que non :
on a toujours pu les éluder. Mais il refuse d'en conclure qu'il faut
continuer, et il refuse l'argument selon lequel la classification d'objets
géométriques « relativement merdiques » se ramènerait à des invariants
cohomologiques bien connus : c'est confondre la \emph{compréhension d'un
objet géométrique} et la détermination de sa classe à équivalence près.

Suivent quatre exemples. Dans le langage des conventions, une
\emph{$2$-catégorie de Picard stricte} est un $2$-groupoïde muni d'une
structure de groupe abélien strictement commutative à cohérence près, et
elle équivaut à la donnée d'un complexe de groupes abéliens concentré en
trois degrés consécutifs, à quasi-isomorphisme près ; ses $\pi_0$, $\pi_1$,
$\pi_2$ sont les trois groupes de cohomologie de ce complexe, du plus haut
degré au plus bas.\footnote{Pour $n = 1$ c'est le théorème de Deligne
(SGA~4, XVIII, 1.4.15). Pour $n$ quelconque, la théorie des
$n$-catégories de Picard strictes n'existait pas en 1975 ; la lettre la
suppose. L'énoncé qu'on donne est le contenu, en termes d'aujourd'hui, de
l'équivalence entre $n$-groupoïdes de Picard stricts et complexes de
longueur $n$ à quasi-isomorphisme près, qu'il demande explicitement page
56 ; on en dit le statut actuel à cet endroit.}

\emph{1) Gerbes.} Soit $L$ un lien sur le topos $X$, de centre $Z$ (un
faisceau abélien). Les gerbes liées par $Z$ forment une $2$-catégorie de
Picard stricte, associée au complexe $\tau_{\leqslant 2}\,R\Gamma(X,Z)$, de
sorte que
\[
\pi_{2-i}\bigl(\mathrm{Gerbes}_Z(X)\bigr) \simeq H^i(X,Z), \qquad 0 \leqslant i \leqslant 2 .
\]
Les gerbes liées par $L$ forment un pseudo-torseur sous cette
$2$-catégorie, non vide si et seulement si une obstruction dans $H^3(X,Z)$
est nulle --- c'est le théorème de Giraud (\emph{Cohomologie non
abélienne}, IV.3.4). Pour comprendre cette classe, il propose de passer aux
champs : le $2$-champ de Picard strict des $Z$-gerbes sur les objets
variables de $X$, qui est le $2$-champ $B^2Z$, et le $2$-champ des
$L$-gerbes, qui est un $2$-torseur sous le précédent ; et les $2$-torseurs
sous $B^2Z$ sont classés par $H^3(X,Z)$, comme les $Z$-gerbes, torseurs sous
le champ de Picard strict $BZ$ des $Z$-torseurs, le sont par
$H^2(X,Z)$.\footnote{La page dit que la $3$-catégorie de Picard stricte des
$2$-gerbes liées par $Z$ est décrite par $R\Gamma(X,Z)$ « tronqué en
dimension 3 », et que le « $3$-champ de Picard correspondant » l'est par une
résolution injective de $Z$ tronquée en degré 3. Pour un champ, ce qui le
détermine est le complexe de faisceaux $Z[3]$ ; une résolution injective
tronquée en est un représentant commode, dont les sections globales
calculent les groupes $H^i(X,Z)$, $0 \leqslant i \leqslant 3$. On écrit
« $B^2Z$ » pour ce qu'il nomme le $2$-champ des $Z$-gerbes, ce qui est
exact parce que toute gerbe est localement triviale. Le préfixe « 1- » de
« $1$-torseurs » est une lecture incertaine. La théorie des $2$-gerbes et de
leur classification par $H^3$ a été écrite depuis par Breen lui-même
(\emph{On the classification of 2-gerbes and 2-stacks}, Astérisque 225,
1994).}

\emph{2) Champs de Picard d'invariants $M$ et $N$.} Si $M$ et $N$ sont deux
faisceaux abéliens, les champs de Picard (non nécessairement stricts) dont
le faisceau des classes d'objets est $M$ et celui des automorphismes de
l'objet nul est $N$ forment une $2$-catégorie de Picard stricte. Il la dit
« représentée sans doute » par le complexe $\mathrm{RHom}(X(M),N)$ tronqué
en degré $2$, d'invariants « le drôle de $\mathrm{Ext}^2$ » de sa lettre à
Deligne et les « honnêtes » $\mathrm{Ext}^0(M,N)$ et
$\mathrm{Ext}^1(M,N)$.\footnote{Le complexe $X(M)$ n'est pas défini sur ces
pages ; l'appel (23), qui le commentait sans doute, renvoie aux notes non
reprises. La lecture la plus sûre est que $\mathrm{RHom}(X(M),N)$ joue le
rôle du complexe $E'(M,N)$ de la lecture modernisée du dossier 134-1 : en
termes d'aujourd'hui, le complexe des morphismes de $M$ dans $N$ calculé au-dessus de
la sphère $\mathbb{S}$ et non de $\mathbb{Z}$, dont les $\mathrm{Ext}^0$ et
$\mathrm{Ext}^1$ sont les ordinaires et dont le $\mathrm{Ext}^2$ classe les
champs de Picard non nécessairement stricts. Cette identification est la
nôtre. Sur la page, l'exposant du premier « $\mathrm{Ext}$ » ressemble à un
1 ; l'encadrement $0 \leqslant i \leqslant 1$ demande $i$, et c'est ce qu'on
lit.} Les champs de Picard \emph{stricts} forment une sous-$2$-catégorie de
Picard pleine, associée à $\tau_{\leqslant 2}\,\mathrm{RHom}(M,N)$ : le
groupe $\mathrm{Ext}^2(M,N)$ classe les objets à équivalence près,
$\mathrm{Ext}^1(M,N)$ les automorphismes de l'objet nul à isomorphisme près,
$\mathrm{Ext}^0(M,N)$ les automorphismes de l'automorphisme identique de
celui-ci. Il n'a pas, dit-il, d'interprétation géométrique de
$\tau_{\leqslant n}\,\mathrm{RHom}(M,N)$ pour $n > 2$, mais il faut la
chercher du côté des $n$-champs de Picard.\footnote{En termes
d'aujourd'hui, et c'est notre commentaire : $\tau_{\leqslant n}\,
\mathrm{RHom}(M,N)$ est la $n$-catégorie de Picard stricte des morphismes
de $HM$ dans $B^{n-1}HN$ dans les faisceaux de $H\mathbb{Z}$-modules, et
c'est bien un objet de la famille qu'il annonce.}

\emph{3) Gr-champs liés par $(G,N)$.} Soit $G$ un faisceau de groupes sur
$X$ opérant sur un faisceau abélien $N$. Les champs en Gr-catégories sur
$X$ liés par $(G,N)$ --- $\pi_0 = G$, $\pi_1 = N$, avec l'action donnée ---
forment une $2$-catégorie de Picard, d'invariants
\[
\pi_0 = H^3(B_G/X, N), \qquad \pi_1 = H^2(B_G/X, N), \qquad \pi_2 = Z^1(G,N),
\]
où $B_G/X$ est le topos classifiant de $G$ au-dessus de $X$ et $Z^1(G,N)$
le groupe des $1$-cocycles (homomorphismes croisés). Il laisse à Breen le
soin de deviner le complexe tronqué qui la décrit.\footnote{Notre réponse,
dans le cas d'un groupe discret $G$ sur un point : le complexe des cochaînes
inhomogènes privé de son degré $0$ et tronqué en degré $3$,
$C^1(G,N) \to C^2(G,N) \to Z^3(G,N)$. Supprimer $C^0$ remplace $H^1$ par
$Z^1$, ce qui est exactement l'invariant $\pi_2$ de la page :
l'automorphisme d'un foncteur monoïdal qui induit l'identité sur $G$ et sur
$N$ est une famille $u_x \in N$ avec $u_{xy} = u_x + x\cdot u_y$, sans
quotient ultérieur. La version faisceautique remplace les cochaînes par
celles du classifiant simplicial de $G$ au-dessus de $X$. La lettre à
Deligne citée ici n'est pas dans le dossier.}

\emph{4) Algèbres d'Azumaya.} Sur un topos localement annelé $X$, les
Algèbres d'Azumaya forment une $2$-catégorie de Picard « de Brauer » : la
catégorie des morphismes $A \to B$ est celle des trivialisations
$(E,\varphi)$ de $A^{\circ}\otimes B$, avec $E$ un Module localement libre
et $\varphi : \underline{\mathrm{End}}(E) \xrightarrow{\sim} A^{\circ}\otimes
B$ --- ce qu'on appelle aujourd'hui une équivalence de Morita ---, le
produit est le produit tensoriel et l'inverse l'algèbre opposée. En
associant à chaque algèbre la gerbe de ses trivialisations, on obtient un
$2$-foncteur monoïdal vers la $2$-catégorie de Picard des $\mathbb{G}_m$-gerbes, qui
est $2$-fidèle (pleinement fidèle sur les catégories de morphismes) ; d'où
les invariants
\[
\pi_0 = H^2(X,\mathbb{G}_m)_{\mathrm{Br}}, \qquad \pi_1 = H^1(X,\mathbb{G}_m) = \mathrm{Pic}(X), \qquad \pi_2 = H^0(X,\mathbb{G}_m),
\]
où $H^2(X,\mathbb{G}_m)_{\mathrm{Br}}$ est le sous-groupe des classes qui
proviennent d'algèbres d'Azumaya. Formées sur les objets variables de $X$,
ces $2$-catégories font une $2$-catégorie de Picard fibrée qui n'est « pas
tout à fait un $2$-champ (whatever that means) » : la condition de
$2$-recollement doit faillir, sans quoi il n'y aurait pas d'indice
$\mathrm{Br}$.\footnote{L'argument est juste et se dit ainsi : le
$2$-champ associé est celui de toutes les $\mathbb{G}_m$-gerbes, puisque
toute gerbe est localement triviale, c'est-à-dire localement celle de
$\mathcal{O}_X$ ; si le préchamp était déjà un $2$-champ, toute classe de
$H^2$ viendrait d'une algèbre. La comparaison du groupe de Brauer et de
$H^2(X,\mathbb{G}_m)_{\mathrm{tors}}$ est sa question de \emph{Le groupe de
Brauer II} ; l'égalité est connue pour les schémas munis d'un faisceau
ample (Gabber, puis de Jong, 2003). Notre commentaire : c'est en admettant
des algèbres d'Azumaya \emph{dérivées} que Toën (\emph{Derived Azumaya
algebras and generators for twisted derived categories}, Invent. Math.,
2012) obtient un objet qui recolle et atteint tout $H^2(X,\mathbb{G}_m)$
pour les schémas quasi compacts et quasi séparés --- ce qui est, à la
lettre, le défaut de $2$-recollement que la page constate, réparé en
élargissant la classe d'objets.}

\pagerange{40}{42}
\subsection*{App.~2 : les $n$-champs et l'algèbre homologique non commutative (pages 40 à 42)}

Les $n$-catégories de Picard strictes, qui « s'imposent pas à pas » dans un
contexte commutatif, donnent une interprétation géométrique correcte des
complexes tronqués à l'ordre $n$ en tant qu'objets de catégories dérivées,
donc, à la limite sur $n$, des complexes tout court. D'où l'idée « naïve »
qu'il propose : les « complexes non commutatifs », objets fantômes d'une
algèbre homologique non commutative, seraient ce qui reste d'une
$n$-catégorie de Picard quand on oublie sa structure de Picard --- la
$n$-catégorie elle-même --- et, sur un topos, les $n$-champs.\footnote{C'est
exactement ce que l'on fait aujourd'hui, et c'est notre commentaire : un
complexe connectif à quasi-isomorphisme près est un $\infty$-groupoïde
muni d'une structure de $H\mathbb{Z}$-module, et oublier cette structure
laisse un objet de $\mathcal{S}$, ou, sur un topos, un faisceau d'espaces.
Le foncteur d'oubli est celui de Dold--Kan suivi de l'oubli de la
structure de groupe simplicial.}

L'idée, dit-il, est d'abord venue d'ailleurs : des théorèmes de Lefschetz
« faibles » en cohomologie étale à coefficients discrets, pour une variété
projective et toute section hyperplane, ou une variété quasi projective et
presque toute section hyperplane, sous les hypothèses de profondeur
cohomologique les plus naturelles. Dans le cas commutatif, la dualité
indique les meilleurs énoncés possibles (exposé de Mme Raynaud dans SGA~2),
mais seulement pour des coefficients premiers aux caractéristiques ; des
démonstrations géométriques directes donnaient des énoncés voisins pour le
$\pi_0$ et le $\pi_1$ sans cette restriction, du moins dans le cas propre.
Les meilleurs énoncés, conjecturés dans SGA~2, ont été démontrés par
Michèle Raynaud dans sa thèse. Le point est que ces énoncés se formulent le
plus naturellement en termes de $1$-champs sur le site étale, la notion de
profondeur $\geqslant i$ ($i = 1,2,3$) aussi, et qu'on ne semble pas pouvoir
démontrer, par exemple, la bijectivité $\pi_1(Y)\to\pi_1(X)$ pour une section
hyperplane $Y$ avec le seul formalisme du $H^0$ et du $H^1$ non abéliens de
Giraud : il faut les trois théorèmes de changement de base --- propre, lisse,
et « propreté cohomologique générique » --- pour les $1$-champs de torsion
eux-mêmes.\footnote{Page 41. Le troisième énoncé dit : pour $f : X \to S$ de
type fini, $S$ intègre, il existe un ouvert non vide $U$ de $S$ tel que la
formule de changement de base vaille pour \emph{tout} changement de base se
factorisant par $U$ (la frappe porte « $u$ »). Pour le changement de base
lisse et le troisième, il faut des faisceaux d'automorphismes premiers aux
caractéristiques, et une hypothèse de constructibilité pour le troisième ;
ils ne servent que dans la version générique du théorème. Il croit, sans en
être sûr, que Giraud démontre les deux premiers, peut-être le troisième,
pour les $1$-champs.} Le contexte naturel des théorèmes de changement de
base et de Lefschetz, et de la profondeur, doit donc être celui des
$n$-champs, et il ne s'agira pas d'une « jonglerie purement formelle » : on
sera aussitôt soumis à des tests d'utilisabilité sérieux.

Il propose enfin d'expliciter les ingrédients d'une telle algèbre, en
rappelant que même l'algèbre homologique commutative n'est pas dans un état
satisfaisant, puisqu'on ne sait pas quelle est « la bonne notion de catégorie
triangulée » : on fait de la prose, comme Monsieur Jourdain, sans se
demander quelles structures portent les catégories de complexes.\footnote{Les
réponses qu'on donne aujourd'hui à cette question sont les catégories
triangulées munies d'un enrichissement différentiel gradué (Bondal--Kapranov,
1990), les $\infty$-catégories stables (Lurie, \emph{Higher Algebra},
chapitre 1) et les dérivateurs, dont \emph{Pursuing Stacks} amorce la
théorie (partie IX du plan de la page 2). Le commentaire est le nôtre ; les
appels (*) et (4) de ce passage sont ajoutés à la main, et un appel marginal
est lu avec doute « (24) ».}

\pagerange{42}{45}
\subsection*{App.~3 : les $n$-groupoïdes comme types d'homotopie tronqués (pages 42 à 45)}

C'est le post-scriptum, écrit pour « mettre en appétit », et c'est le cœur
des trois lettres. Le « yoga » est le suivant : une $n$-catégorie en
groupoïdes, à $n$-équivalence près, « est essentiellement la même chose »
qu'un ensemble simplicial à homotopie près dont on néglige les $\pi_i$ pour
$i \geqslant n+1$.

\emph{Le $n$-groupoïde fondamental d'un complexe de Kan.} À un ensemble
simplicial de Kan $K_\bullet$, il associe une $n$-catégorie
$C_n(K_\bullet)$ : ses $0$-objets sont les $0$-simplexes, ses $1$-flèches les
chemins, ses $2$-flèches les homotopies entre chemins à extrémités fixes, et
ainsi de suite jusqu'aux $n$-flèches, qui sont les homotopies prises
\emph{à homotopie près}. La composition est associative à homotopie près
seulement, de sorte que l'objet obtenu n'est pas strict ; et c'est
précisément la mise au point de ce yoga qui pourrait servir de fil d'Ariane
pour définir les $n$-catégories non strictes, leurs $n$-foncteurs et leurs
$n$-équivalences, à côté du yoga inductif « une $n$-catégorie est une
catégorie dont les $\mathrm{Hom}$ sont des $(n-1)$-catégories ». La
construction est fonctorielle, et même $n$-fonctorielle en $K_\bullet$.\footnote{Notre commentaire sur le caractère non strict, qui est le point
technique. Les $n$-groupoïdes \emph{stricts} ne suffisent pas : Simpson
(\emph{Homotopy types of strict 3-groupoids}, 1998) a montré qu'aucun
$3$-groupoïde strict n'a le $3$-type de $S^2$, ce qui invalide l'énoncé
contraire de Kapranov et Voevodsky (1991). Pour des modèles non stricts,
l'équivalence entre $n$-groupoïdes et $n$-types est un théorème dans
plusieurs cadres --- les $n$-groupoïdes de Tamsamani (1999), les complexes
de Kan $n$-tronqués, où elle est presque tautologique --- et elle est
restée longtemps ouverte pour les $\infty$-groupoïdes globulaires que
Grothendieck définira dans \emph{Pursuing Stacks} ; nous ne garantissons
pas l'état actuel de ce dernier cas. Page 43, la frappe porte « $K!$ » pour
$K'_\bullet$, et la phrase sur la $n$-catégorie des ensembles simpliciaux
se termine par « quel que soit $N$ ».}

\emph{Les groupes d'homotopie.} $C_n(K_\bullet)$ est un $n$-groupoïde : toute
$i$-flèche, $1 \leqslant i \leqslant n$, y est une équivalence. Pour un
$n$-groupoïde $C$ et un objet $x$, on pose : $\pi_0(C)$ l'ensemble des
classes d'objets ; $\pi_1(C,x)$ le groupe des classes de $1$-flèches
$x \to x$ ; $\pi_2(C,x)$ celui des classes de $2$-flèches
$1_x \Rightarrow 1_x$, qui est commutatif, comme les suivants ; et ainsi de
suite jusqu'à $\pi_n$. Ce sont des systèmes locaux sur $\pi_0$, et il y a
des isomorphismes canoniques $\pi_i(K_\bullet,x) \simeq \pi_i(C_n(K_\bullet),x)$
pour $0 \leqslant i \leqslant n$. Par suite, \emph{une application
simpliciale $f : K_\bullet \to K'_\bullet$ induit une $n$-équivalence
$C_n(K_\bullet) \to C_n(K'_\bullet)$ si et seulement si elle induit une
bijection sur les $\pi_i$ pour $0 \leqslant i \leqslant n$, en tout point
base.}

Il hésite alors : il « serait plus heureux » de pouvoir ajouter « et une
surjection sur $\pi_{n+1}$ », condition qui lui paraît nécessaire pour
identifier la localisée des ensembles simpliciaux selon ces flèches à celle
des $n$-catégories selon les $n$-équivalences.\footnote{L'hésitation n'a pas
lieu d'être, et c'est notre correction. Les applications qui induisent une
bijection sur les $\pi_i$, $i \leqslant n$, sont exactement celles que la
troncature de Postnikov $\tau_{\leqslant n}$ rend inversibles ; la localisée
des ensembles simpliciaux selon elles est équivalente à la catégorie
homotopique des $n$-types, puisque $K_\bullet \to \tau_{\leqslant n}K_\bullet$
est l'une d'elles et qu'entre $n$-types elles sont des équivalences faibles.
Elle est donc équivalente à celle des $n$-groupoïdes dès qu'on dispose de
l'hypothèse d'homotopie. La condition de surjectivité sur $\pi_{n+1}$ définit
autre chose : les applications $(n+1)$-connexes, qui reviendront page 48, où
le même décalage d'un degré se produit. Page 44, « $i = n+1$ » est frappé
« $N+1$ », le $N$ surchargé en $n$.} Il ajoute que la troncature homotopique
(tuer les $\pi_i$ pour $i > n$) se traduit par l'opération qui conserve les
$i$-flèches pour $i < n$ et remplace les $n$-flèches par leurs classes à
équivalence près, et l'inclusion des $n$-types dans les $N$-types par l'ajout
d'identités en dimensions $n+1$ à $N$ --- ce qui est exact, la troncature
étant adjointe à gauche de l'inclusion.

\emph{Les $\infty$-catégories.} Sauf erreur, dit-il, la localisée des
$\infty$-catégories (en groupoïdes) selon les $\infty$-équivalences est la
catégorie homotopique, « ou du moins une sorte de complétée de
celle-là ».\footnote{Notre commentaire : pour les espaces, la tour de
Postnikov converge et il n'y a pas de complétion à faire ; la réserve devient
pertinente dans un topos, où un faisceau d'espaces n'est pas toujours limite
de ses tronqués --- c'est la distinction entre un $\infty$-topos et son
hypercomplété (Lurie, \emph{Higher Topos Theory}, § 6.5).} Une
$\infty$-catégorie de Picard stricte s'interprète comme un complexe de
chaînes regardé dans une catégorie dérivée, ce qu'il faut relier à la
théorie de Dold--Puppe, qui interprète ces complexes comme des groupes
abéliens simpliciaux.\footnote{L'énoncé strict correspondant --- les
groupes abéliens dans les $\omega$-groupoïdes stricts sont les complexes de
chaînes en degrés positifs --- est dû, à notre connaissance, à Brown et
Higgins (1981) et Bourn (1990).}

\emph{Le dictionnaire.} Pour se donner confiance, il traduit des
constructions familières. L'espace des lacets $\Omega(K_\bullet,x)$
correspond à la $(n-1)$-catégorie $\underline{\mathrm{Hom}}(x,x)$. L'espace
simplicial des applications $\mathrm{Hom}_\bullet(K_\bullet,K'_\bullet)$
correspond à la $n$-catégorie $\underline{\mathrm{Hom}}(C,C')$. La fibre
homotopique d'une application correspond au produit fibré homotopique
$C \times_{C'} \mathrm{pt}$ dans les $n$-groupoïdes.\footnote{La page dit
« produit $n$-fibré », puis une main corrige en « $(n+1)$ » sans que la
portée de la correction soit sûre. La correction est juste : le produit
fibré homotopique de $n$-catégories est une limite dans la
$(n+1)$-catégorie qu'elles forment, comme le « produit $2$-fibré » de
catégories de Giraud est celui des $1$-catégories.} Les espaces
$K(\pi,n)$ sont les $n$-gerbes liées par $\pi$. Il ne voit pas de candidat
pour la suspension.\footnote{Il n'y en a pas à l'intérieur des
$n$-groupoïdes : la suspension d'un $n$-type n'est pas un $n$-type. Ce qui
existe est $\tau_{\leqslant n}\Sigma$, la somme amalgamée
$\mathrm{pt}\sqcup_C\mathrm{pt}$ calculée dans les $n$-groupoïdes.
Commentaire nôtre.} Enfin le dévissage de Postnikov : un $n$-groupoïde
s'interprète comme son tronqué $(n-1)$-groupoïde muni d'une $n$-gerbe sur
le topos classifiant de celui-ci, liée par $\pi_n$ tordu par l'action de
$\pi_1$. C'est la description actuelle de l'étage
$\tau_{\leqslant n}K \to \tau_{\leqslant n-1}K$ comme fibration de fibre
$K(\pi_n,n)$, classée par un invariant de Postnikov
$k_n \in H^{n+1}(\tau_{\leqslant n-1}K;\pi_n)$ à coefficients locaux.

\pagerange{45}{46}
\subsection*{App.~4 : relativisation sur un topos (pages 45 et 46)}

Tout ce yoga doit être relativisé au-dessus d'un topos $X$ : il faudrait
identifier, dans une certaine mesure, l'algèbre homotopique des faisceaux
simpliciaux sur $X$ et l'algèbre « catégorique » des $n$-champs en
groupoïdes, $0 \leqslant n \leqslant \infty$.\footnote{C'est ce que
réalise aujourd'hui la structure de modèles de Jardine sur les faisceaux
simpliciaux (1987), dont la localisée est l'hypercomplété de
l'$\infty$-topos $\mathrm{Sh}(X;\mathcal{S})$ (Lurie, \emph{Higher Topos
Theory}, § 6.5.2). Les faisceaux simpliciaux d'Illusie et la théorie
homotopique abstraite de K.~Brown (1973) sont antérieurs à la lettre, que ni
l'un ni l'autre ne cite. Commentaire nôtre.} L'image inverse des faisceaux
simpliciaux, qui est évidente, doit correspondre à l'image inverse plus
subtile des $n$-champs ; et l'image directe des $n$-champs, qui est évidente,
à une image directe dérivée des faisceaux simpliciaux --- « on hésite s'il
faut mettre $Lf_*$ ou $Rf_*$ ».\footnote{C'est $Rf_*$ : $f_*$ est adjoint
à droite de $f^*$, et on le dérive à droite, en remplaçant le faisceau
simplicial par un remplacement fibrant (localement de Kan et satisfaisant
la descente). Sur la page, « $Lf_*(K_\bullet)$ » est écrit au-dessus de
« d'un faisceau », qu'un trait semble barrer en partie.} Les dévissages de
Postnikov doivent avoir une interprétation simple en termes de champs ; il
rappelle la remarque de Giraud qu'un $1$-champ en groupoïdes équivaut à la
donnée de son faisceau $\pi_0$ et d'une gerbe sur le topos induit
$X_{/\pi_0}$, gerbe dont le lien « devrait être noté $\pi_1$ ». Cela vaut
pour tout $n$ : un $n$-champ $F$ est un objet connexe de l'$\infty$-topos
$\mathrm{Sh}(X;\mathcal{S})_{/\pi_0 F}$.\footnote{Notre commentaire. La
traduction des $1$-champs en topos classifiants au-dessus de $X$ est, dit-il,
traitée par Giraud sous le nom d'« extensions » du topos. Après « Giraud »
la frappe porte un « x » en indice, peut-être un appel de note.}

\emph{Remords.} Mais il ne sait pas définir à vue de nez le topos classifiant
d'un $n$-champ pour $n \geqslant 2$ --- il suffirait de savoir le faire pour
une $n$-gerbe liée par un faisceau abélien $\pi_n$ --- et ne sait donc pas
décrire le dévissage de Postnikov en termes de $n$-champs au-delà de
$n \leqslant 2$. C'est lié à la définition directe de la cohomologie d'un
$n$-groupoïde, qui devrait redonner celle de l'ensemble simplicial
correspondant : peut-on le faire par un nerf ? Cas typique : trouver, pour un
groupe commutatif $\pi$, un topos canonique et fonctoriel en $\pi$ qui ait
le type d'homotopie de $K(\pi,n)$. « On frémit à l'idée que les topos
pourraient ne pas faire l'affaire, et qu'il y faille des “$n$-topos” !!
(J'espère bien que ces animaux n'existent pas …) »\footnote{Notre
commentaire. Il y a deux questions, et elles ont deux réponses. \emph{Le
type d'homotopie} : des topos ordinaires suffisent ; le topos des
préfaisceaux sur la catégorie des simplexes d'un ensemble simplicial $K$ a
le type d'homotopie faible de $K$, ce qui donne, avec un modèle simplicial
fonctoriel de $K(\pi,n)$, le topos canonique demandé (voir Moerdijk,
\emph{Classifying spaces and classifying topoi}, LNM 1616, 1995, et
Illusie). C'est d'ailleurs le principe des « catégories test » de
\emph{Pursuing Stacks}. \emph{La propriété universelle} : pour qu'un objet
classifie les $n$-gerbes comme un topos classifiant classe les torseurs, il
faut des $\infty$-topos (Toën--Vezzosi, Rezk, Lurie) ; les $n$-topos
existent (Lurie, \emph{Higher Topos Theory}, § 6.4), et ils sont devenus le
cadre ordinaire de ce qu'il cherchait. On y revient page 50.} Pour un
$n$-groupoïde $C$, on dispose bien d'un nerf, et la cohomologie de $C$ à
coefficients dans un système local est celle de ce nerf.

\pagerange{46}{47}
\subsection*{App.~5 : une « algèbre topologique » (pages 46 et 47)}

L'« algèbre homologique non commutative » qu'il suggère serait l'étude
parallèle de quatre familles de notions et de leurs relations : a) espaces
topologiques et topos ; b) ensembles et faisceaux simpliciaux ; c)
$n$-catégories, notamment $n$-groupoïdes, et $n$-champs ; d) complexes de
groupes et de faisceaux abéliens. C'est « de l'algèbre avec la présence
constante de motivations provenant de l'intuition topologique ». Il lui
propose le nom d'« algèbre topologique », l'usage courant de ce terme
n'étant guère qu'un « bric à brac » d'anneaux, de corps et de groupes
topologiques, et rappelle que Serre a bien donné un sens nouveau à
« géométrie analytique » sans rencontrer de résistance.\footnote{Le nom n'a
pas pris. Ce domaine s'appelle aujourd'hui, selon l'angle, algèbre
homotopique, théorie des catégories supérieures ou théorie des
$\infty$-topos. La correction manuscrite de la phrase de la page 47 est lue
avec doute.} La lettre se termine par « Re-salut, et au plaisir de te
lire. A.G. »

\pagerange{48}{55}
\section*{Deuxième lettre à Breen, Villecun, 17 février 1975 (pages 48 à 55)}

La lettre, en français, est une copie dactylographiée par Breen, dont
l'intitulé « Transcription d'une lettre manuscrite de A. Grothendieck » est
barré ; les lettres cursives y sont tracées à la main. Elle commence : « Encore
un “afterthought” à une lettre-fleuve sur le yoga homotopique ».

\pagerange{48}{51}
\subsection*{App.~6 : le $n$-groupoïde fondamental d'un topos (pages 48 à 51)}

\emph{Le pro-type d'homotopie.} À un topos $X$ on associe canoniquement un
pro-ensemble simplicial, donc un pro-type d'homotopie (Artin--Mazur,
\emph{Etale homotopy}, LNM 100, 1969). Si $X$ est « localement
homotopiquement trivial », ce pro-objet est essentiellement constant et $X$
a un type d'homotopie ordinaire ; s'il l'est en dimension $\leqslant n$, il
a un type d'homotopie $n$-tronqué.\footnote{Ce pro-type d'homotopie est
aujourd'hui la \emph{forme} (shape) de l'$\infty$-topos associé (Lurie,
\emph{Higher Topos Theory}, § 7.1.6), et « localement homotopiquement
trivial » correspond à « localement de forme constante ».}

\emph{Trois critères, et leur décalage.} Pour un morphisme de topos
$f : X \to Y$, il cite le critère cohomologique d'Artin--Mazur : $f$ est une
« équivalence d'homotopie en dim $\leqslant n$ » si $H^i(Y,F)\to
H^i(X,f^*F)$ est bijectif pour $i \leqslant n$ et tout faisceau de groupes
localement constant $F$, avec $i \leqslant 1$ si $F$ n'est pas
commutatif. Il le traduit en termes de $n$-gerbes localement constantes,
puis affirme qu'il équivaut encore à chacun des deux critères suivants :

\begin{quote}
(A) pour tout $n$-champ localement constant $\mathcal{F}$ sur $Y$, le
$n$-foncteur $\mathcal{F}(Y) \to (f^*\mathcal{F})(X)$ est une
$n$-équivalence ;

(B) le $n$-foncteur $\mathcal{F}\mapsto f^*\mathcal{F}$, de la $n$-catégorie
des $(n-1)$-champs localement constants sur $Y$ dans celle des
$(n-1)$-champs localement constants sur $X$, est une $n$-équivalence.
\end{quote}

Les trois énoncés sont justes chacun pour lui-même, mais ils ne sont pas
équivalents : ils diffèrent d'un degré. Exactement, pour $X$ et $Y$
localement de forme constante et $f$ bijectif sur $\pi_0$ et
$\pi_1$ :\footnote{Correction nôtre ; la page affirme l'équivalence des
trois (« il est certainement vrai que ceci équivaut encore … »). La flèche
$\mathcal{F}(Y)\to\mathcal{F}(X)$ de la page est tracée à la main, avec une
boucle à l'origine, peut-être un $\sim$, et le critère par les gerbes porte
la clause « pour tout tel $\mathcal{F}$ et $i \leqslant n$ », qu'on lit
comme on le dit ici. Les deux derniers points se déduisent de ce que les
applications $(n+1)$-connexes sont exactement celles qui sont orthogonales
à gauche aux applications $n$-tronquées : les sections d'un système local à
fibres $n$-tronquées se relèvent alors de façon unique. Les conditions sur
les coefficients doivent s'entendre à coefficients locaux.}
\begin{itemize}
  \item le critère cohomologique sous la forme citée (bijection sur $H^i$
  pour $i \leqslant n$) caractérise les $f$ \emph{$n$-connexes} : bijection
  sur $\pi_i$ pour $i < n$, surjection sur $\pi_n$ ;
  \item le critère par les $n$-gerbes, \emph{à condition d'y inclure les
  gerbes non neutres}, et le critère (A) caractérisent tous deux les $f$
  \emph{$(n+1)$-connexes} : bijection sur $\pi_i$ pour $i \leqslant n$,
  surjection sur $\pi_{n+1}$ ; cohomologiquement, bijection sur $H^i$ pour
  $i \leqslant n$ et \emph{injection} sur $H^{n+1}$ --- c'est cette
  injectivité qu'exprime le fait qu'une $n$-gerbe sur $Y$ dont l'image
  inverse a une section en a déjà une ;
  \item le critère (B) caractérise exactement les $f$ qui induisent une
  bijection sur $\pi_i$ pour $i \leqslant n$, c'est-à-dire tels que
  $\tau_{\leqslant n}f$ soit une équivalence : la pleine fidélité de $f^*$
  est la $n$-connexité, l'essentielle surjectivité ajoute l'injectivité sur
  $\pi_n$.
\end{itemize}
On a donc (A) $\Rightarrow$ (B) $\Rightarrow$ critère cohomologique cité,
et aucune réciproque. Deux exemples le montrent. L'inclusion d'un point
dans un espace $K(A,n+1)$, $A \neq 0$, vérifie (B) mais non (A) : la
fibration des chemins, de fibre $K(A,n)$, est un $n$-champ localement
constant qui a des sections sur le point et n'en a pas sur $K(A,n+1)$.
L'inclusion $S^2\vee S^2 \to S^2\times S^2$ induit des bijections sur les
$H^i$, $i \leqslant 3$, pour tout groupe de coefficients, mais non sur
$\pi_3$ (qui est $\mathbb{Z}^3$ d'un côté, $\mathbb{Z}^2$ de l'autre) : elle
vérifie le critère cité pour $n = 3$, mais ni (B) ni (A).

Ce que la lettre veut en tirer reste vrai, et c'est (B) qui le porte : les
constructions qu'on peut faire sur $X$ en termes de $(n-1)$-champs
localement constants ne dépendent que de son type d'homotopie
$n$-tronqué, et le déterminent.

\emph{Le $n$-groupoïde fondamental.} Lorsque $X$ est localement
homotopiquement trivial en dimension $\leqslant n$, son type d'homotopie
$n$-tronqué s'interprète comme un $n$-groupoïde, défini à $n$-équivalence
près, qu'il propose d'appeler le \emph{$n$-groupoïde fondamental} de $X$ et
de noter $\Pi_n(X)$ ; et

\begin{quote}
(C) les $(n-1)$-champs localement constants sur $X$ s'identifient aux
$n$-foncteurs de $\Pi_n(X)$ dans la $n$-catégorie $((n-1)\text{-}\mathrm{Cat})$
de toutes les $(n-1)$-catégories.
\end{quote}

Pour $n = 1$ c'est la classification des revêtements par le
groupoïde fondamental, « la théorie de Poincaré ». $\Pi_n(X)$ détermine les
$\pi_i(X)$, $i \leqslant n$, et les invariants de Postnikov jusqu'à celui
qui vit dans $H^{n+1}(\Pi_{n-1}(X),\pi_n)$. Pour un topos quelconque, il
espère un pro-$n$-groupoïde ; c'est fait pour $n = 1$ et $X$ connexe, et
traité « in extenso » dans SGA~3 pour le topos étale d'un schéma, à propos de
la classification des tores.\footnote{Notre commentaire. L'énoncé (C) est
aujourd'hui un théorème : pour un $\infty$-topos localement de forme
constante, les faisceaux localement constants à valeurs dans une
$\infty$-catégorie $\mathcal{D}$ forment l'$\infty$-catégorie
$\mathrm{Fun}(\Pi_\infty(X),\mathcal{D})$ (Lurie, \emph{Higher Algebra},
appendice A) ; si $\mathcal{D}$ est formée de $(n-1)$-catégories, un tel
foncteur ne voit que $\Pi_n(X)$, puisqu'il se factorise par le
sous-groupoïde maximal de $\mathcal{D}$, qui est $n$-tronqué. Le cas d'un
$\infty$-topos quelconque, avec un pro-groupoïde, est la « théorie de Galois
supérieure » de Hoyois (J. Pure Appl. Algebra, 2018). Sur la page, un
crochet hachuré marqué « (C) » encadre l'énoncé. La lettre écrit ce
paragraphe avec $(n-1)$-champs « localement constants », et
$((n-1)\text{-}\mathrm{Cat})$ comme but ; c'est la forme juste, et on la
garde.}

\emph{Le $n$-topos classifiant.} Pour $n = 1$, on récupère le groupoïde
$\mathcal{C}_1 = \Pi_1(X)$ à partir du topos « multigaloisien »
$\underline{\mathrm{Hom}}(\mathcal{C}_1,(\mathrm{Ens}))$ des systèmes locaux,
qui est le topos classifiant de $\mathcal{C}_1$, comme la catégorie de ses
foncteurs fibres, c'est-à-dire de ses points.\footnote{La page dit « la
catégorie opposée à la catégorie des points de ce topos » ; pour un groupoïde
la différence est sans objet, l'inversion des flèches identifiant un
groupoïde à son opposé.} Pour élucider (B), il faudrait de même récupérer
$\mathcal{C}_n$ à partir de la $n$-catégorie
\[
\mathcal{T}_n = \underline{\mathrm{Hom}}_n\bigl(\mathcal{C}_n,((n-1)\text{-}\mathrm{Cat})\bigr)
\]
des $(n-1)$-systèmes locaux : $\mathcal{C}_n$ devrait être la
$n$-catégorie des « $n$-foncteurs fibres » sur $\mathcal{T}_n$, ayant des
propriétés d'exactitude analogues à celles d'une image inverse de topos
(commuter aux limites inductives quelconques et aux limites projectives
finies). « C'est ici que se matérialise la peur […] qu'on finisse par tomber
sur la notion de $n$-topos et morphismes de tels ! » --- $\mathcal{T}_n$
serait un $n$-topos, $((n-1)\text{-}\mathrm{Cat})$ le $n$-topos ponctuel, et
$\mathcal{C}_n$ la $n$-catégorie des $n$-points de $\mathcal{T}_n$.
« Brr ! »\footnote{Les cursives $\mathcal{C}_n$ et $\mathcal{T}_n$ se
ressemblent sur la page et ont été départagées par le sens. Notre
commentaire : c'est exactement ce qui est vrai. Pour un $\infty$-groupoïde
$C$, l'$\infty$-topos $\mathcal{S}_{/C} \simeq \mathrm{Fun}(C,\mathcal{S})$
est son $\infty$-topos classifiant, et ses points --- les foncteurs exacts
à gauche qui commutent aux colimites vers $\mathcal{S}$ --- forment un
$\infty$-groupoïde équivalent à $C$, puisqu'un ind-objet d'un groupoïde est
essentiellement constant. Sa version $n$-tronquée est un $n$-topos au sens
de Lurie.}

Peut-on se contenter d'un topos ordinaire ? Le seul problème universel qu'il
voie est celui-ci : trouver un topos $\mathcal{B}$ et un
$\varphi : \Pi_n(\mathcal{B}) \to \mathcal{C}_n$ tels que, pour tout topos
$T$, le foncteur
\[
\underline{\mathrm{Hom}}_{(\mathrm{Top})}(T,\mathcal{B}) \longrightarrow
\tau_1\,\underline{\mathrm{Hom}}\bigl(\Pi_n(T),\mathcal{C}_n\bigr), \qquad
u \longmapsto \varphi\circ\Pi_n(u),
\]
soit une équivalence. Pour $n = 1$, le topos classifiant usuel convient ;
pour $n = 2$, il doute que le problème ait une solution.\footnote{Sur la
page, l'étiquette manuscrite de la flèche porte « $\pi_n(u)$ » pour
$\Pi_n(u)$ ; une main a écrit en marge, en face de ces lignes, « cela semble
suspect ! ». Notre commentaire : la version qui est vraie remplace le topos
ordinaire par l'$\infty$-topos $\mathcal{S}_{/\mathcal{C}_n}$ et la
troncature $\tau_1$ par l'$\infty$-groupoïde tout entier : les morphismes
géométriques $T \to \mathcal{S}_{/C}$ sont les sections globales sur $T$ du
faisceau constant de valeur $C$, c'est-à-dire, par définition de la forme,
les applications de $\Pi_\infty(T)$ dans $C$. Nous ne tranchons pas si le
problème posé avec des topos ordinaires et $\tau_1$ a une solution pour
$n = 2$.}

Il rattache ce doute à un autre : le théorème de van Kampen, qui dit que
$T \mapsto \Pi_1(T)$ transforme sommes amalgamées en sommes amalgamées, ne
serait « sans doute plus vrai » pour $\Pi_2$. Si $T$ est réunion de deux
fermés $T_1$ et $T_2$, il ne serait plus vrai qu'un $1$-champ localement
constant sur $T$ équivaille à la donnée de tels champs sur $T_1$ et $T_2$
et d'une équivalence entre leurs restrictions à $T_1\cap T_2$, alors que
l'énoncé analogue pour les revêtements est correct.\footnote{Correction
nôtre : l'énoncé de recollement est vrai, et van Kampen vaut pour $\Pi_n$
à condition de prendre les sommes amalgamées au sens homotopique. Pour un
recouvrement par deux ouverts, c'est simplement la descente des champs (il
n'y a pas de condition de cocycle non triviale avec deux ouverts) ; pour
deux fermés, il faut une hypothèse de régularité --- par exemple deux
sous-complexes d'un CW-complexe, ou des fermés admettant des voisinages qui
se rétractent sur eux par déformation. Sous ces hypothèses, la forme de $T$
est la somme amalgamée homotopique de celles de $T_1$ et $T_2$ au-dessus de
$T_1\cap T_2$ (Lurie, \emph{Higher Algebra}, A.3.1, pour les ouverts), et
$\tau_{\leqslant n}$, adjoint à gauche, la préserve : $\Pi_n$ commute aux
sommes amalgamées homotopiques, ce qui était déjà le contenu du van Kampen
classique pour $\Pi_1$ dès qu'on le lit $2$-catégoriquement, comme le fait la
page (« commute aux $2$-limites inductives »). Le théorème de van Kampen
supérieur de Brown et Higgins (1981) en est une version stricte. Sur la
page, la phrase est marquée d'un crochet en marge.}

\pagerange{51}{54}
\subsection*{App.~7 : la cohomologie comme intégration de champs (pages 51 à 54)}

\emph{L'intégration.} Si $\mathcal{C}_n = \Pi_n(X)$, $\mathcal{F}$ un
$(n-1)$-champ localement constant sur $X$ et $e^X_{n-1}$ le $(n-1)$-champ
final, on a
\[
\Gamma_X(\mathcal{F}) = \mathcal{F}(X) \simeq \underline{\mathrm{Hom}}(e^X_{n-1},\mathcal{F}),
\]
et, via (C), « l'intégration sur $X$ » des champs localement constants, qui
contient la cohomologie non commutative localement constante de $X$ en
dimension $\leqslant n-1$, devient
$\underline{\mathrm{Hom}}(e^{\mathcal{C}_n}_{n-1},\mathcal{F})$, où
$\mathcal{F}$ est vu comme un $n$-foncteur
$\mathcal{C}_n \to ((n-1)\text{-}\mathrm{Cat})$ et $e^{\mathcal{C}_n}_{n-1}$
est le $n$-foncteur constant de valeur la $(n-1)$-catégorie finale. C'est,
en termes d'aujourd'hui, la limite du diagramme $\mathcal{F}$ indexé par
$\Pi_n(X)$.

\emph{Le cas de Picard.} Soit $\mathcal{F}$ un $n$-champ de Picard strict sur
$X$, défini par un complexe de cochaînes de faisceaux abéliens
$\mathcal{L}^\bullet$ concentré en degrés $0$ à $n$, à quasi-isomorphisme
près. Alors
\[
H^i(X,\mathcal{L}^\bullet) \simeq \pi_{n-i}\,\Gamma_X(\mathcal{F}), \qquad 0 \leqslant i \leqslant n,
\]
le premier membre étant l'hypercohomologie.\footnote{L'indexation se
vérifie ainsi : $\mathcal{F}$ correspond au complexe $\mathcal{L}^\bullet[n]$,
concentré en degrés $-n$ à $0$, et
$\pi_k\Gamma_X(\mathcal{F}) = H^{-k}\bigl(R\Gamma(X,\mathcal{L}^\bullet[n])\bigr)
= H^{n-k}(X,\mathcal{L}^\bullet)$.} Pour atteindre les $H^i$ avec
$i > n$, on regarde $\mathcal{L}^\bullet$ comme concentré en degrés $0$ à
$N$, ce qui remplace $\mathcal{F}$ par $B^{N-n}\mathcal{F}$, où $B$ est le
foncteur espace classifiant --- sur une $n$-catégorie de Picard stricte, il
décale les $i$-flèches en $(i+1)$-flèches et ajoute un unique objet. D'où
la définition, pour tout $n$-champ de Picard strict $\mathcal{F}$,
\[
\boxed{\;H^i(X,\mathcal{F}) = \pi_{N-i}\,\Gamma_X\bigl(B^{N-n}\mathcal{F}\bigr), \qquad N \geqslant \sup(i,n),\;}
\]
indépendante de $N$ grâce au morphisme canonique
$B\,\Gamma_X(\mathcal{F}) \to \Gamma_X(B\mathcal{F})$, qui induit des
isomorphismes sur les $\pi_j$ pour $1 \leqslant j \leqslant
n+1$.\footnote{La page appelle ce morphisme un « morphisme canonique de
$(n-1)$-groupoïdes » ; ce sont des $(n+1)$-groupoïdes, $B\mathcal{F}$ étant
un $(n+1)$-champ. Il n'est pas un isomorphisme sur $\pi_0$ : le membre de
gauche est connexe, et $\pi_0\Gamma_X(B\mathcal{F}) = H^{n+1}(X,\mathcal{F})$.
La formule encadrée l'est sur la page ; on y a écrit $B$ pour sa cursive,
voir les conventions.}

\emph{Sans structure de Picard.} Pour un $n$-champ en groupoïdes quelconque,
on peut encore poser $H^i(X,\mathcal{F}) = \pi_{n-i}\Gamma_X(\mathcal{F})$
pour $0 \leqslant i \leqslant n$, pourvu qu'on se donne une section globale
comme point base pour $i < n$ ; pour un $n$-Gr-champ, le champ classifiant
$B\mathcal{F}$ existe et donne en outre
\[
H^{n+1}(X,\mathcal{F}) = \pi_0\Gamma_X(B\mathcal{F}) = \text{sections de } B\mathcal{F} \text{ à équivalence près},
\]
généralisation directe du $H^1$ non abélien de Giraud, classes de
torseurs.\footnote{Le choix d'une section base est implicite sur la page ;
sans elle, $\pi_{n-i}$ n'a pas de sens pour $n-i \geqslant 1$.} Mais pour
former $B^2\mathcal{F}$ et $H^{n+2}$, il faut que $B\mathcal{F}$ soit
lui-même un Gr-champ, « ce qui ne sera sans doute le cas que si
$\mathcal{F}$ est un $n$-champ de Picard strict ».\footnote{Correction
nôtre : il faut et il suffit que la loi de groupe de $\mathcal{F}$ soit
tressée (structure $E_2$) pour que $B^2\mathcal{F}$ existe, et qu'elle soit
symétrique à cohérence près (structure $E_\infty$, champ de Picard non
nécessairement strict) pour que tous les $B^k\mathcal{F}$ existent. La
stricte commutativité n'est pas nécessaire ; elle caractérise ceux qui
proviennent d'un complexe de faisceaux abéliens. C'est la distinction même
que la lettre à Deligne du dossier 134-1 mesure par un signe.}

\emph{Le cas localement constant.} Si $\mathcal{F}$ est un $n$-champ
localement constant, défini par un $(n+1)$-foncteur
$\mathcal{C}_{n+1} \to (n\text{-catégories de Picard strictes})$, on pose
$H^i(\mathcal{C}_{n+1},\mathcal{F}) = \pi_{n-i}\,
\underline{\mathrm{Hom}}(e^{\mathcal{C}_{n+1}}_n,\mathcal{F})$, et on a
$H^i(\mathcal{C}_{n+1},\mathcal{F}) \simeq H^i(X,\mathcal{F})$ pour
$0 \leqslant i \leqslant n$, sans structure de Picard. Plus généralement,
pour un $\infty$-groupoïde $\mathcal{C}$ et un tel foncteur,
\[
H^i(\mathcal{C},\mathcal{F}) = \pi_{N-i}\,\underline{\mathrm{Hom}}\bigl(e^{\mathcal{C}}_N, B^{N-n}\mathcal{F}\bigr), \qquad N \geqslant \sup(i,n),
\]
c'est-à-dire la cohomologie de $\mathcal{C}$ à coefficients locaux dans un
complexe ; et, pour une simple Gr-structure,
$H^{n+1}(\mathcal{C},\mathcal{F}) = \pi_0\,\underline{\mathrm{Hom}}(e^{\mathcal{C}}_{n+1},B\mathcal{F})$.
Il pense que, pour $\mathcal{C} = \Pi_{n+1}(X)$, cet ensemble doit être
encore canoniquement isomorphe à $H^{n+1}(X,\mathcal{F})$, le vérifie pour
$n = 0$, et demande de décrire la flèche.\footnote{C'est vrai, et c'est
notre réponse. $B\mathcal{F}$ est un système local à fibres
$(n+1)$-tronquées défini sur $\Pi_{n+1}(X)$ ; son espace total au-dessus de
$\Pi_{n+1}(X)$ est donc $(n+1)$-tronqué, et ses sections sur $X$ sont les
mêmes que sur $\Pi_{n+1}(X)$, par la propriété universelle de la troncature.
La flèche est l'image inverse le long du morphisme canonique de $X$ (de sa
forme) vers $\Pi_{n+1}(X)$. Le cas $n = 0$ est la classification des
torseurs sous un groupe localement constant par les revêtements.}

\emph{(A$'$) et (B$'$).} Pour un $(n+1)$-foncteur
$f_{n+1} : \mathcal{C}_{n+1} \to \mathcal{D}_{n+1}$ entre
$(n+1)$-groupoïdes, de tronqué $f_n : \mathcal{C}_n \to \mathcal{D}_n$, la
traduction de (A) et (B) qu'il propose est :

\begin{quote}
(A$'$) $f_n$ est une $n$-équivalence si et seulement si
$f_n^* : \underline{\mathrm{Hom}}(\mathcal{D}_n,((n-1)\text{-}\mathrm{Cat})) \to
\underline{\mathrm{Hom}}(\mathcal{C}_n,((n-1)\text{-}\mathrm{Cat}))$ est une
$n$-équivalence ;

(B$'$) $f_n$ est une $n$-équivalence si et seulement si, pour tout
$n$-système local $\mathcal{F} : \mathcal{D}_{n+1} \to (n\text{-}\mathrm{Cat})$,
le $n$-foncteur
$\Gamma_{\mathcal{D}_{n+1}}(\mathcal{F}) \to \Gamma_{\mathcal{C}_{n+1}}(f_{n+1}^*\mathcal{F})$
est une $n$-équivalence.
\end{quote}

(A$'$) est exact ; c'est la forme combinatoire de (B), et de la
théorie de Galois. (B$'$) ne l'est que dans un sens : la condition sur les
sections équivaut à ce que $f_{n+1}$ soit $(n+1)$-connexe, ce qui demande, en
plus de l'équivalence $f_n$, la surjectivité sur
$\pi_{n+1}$.\footnote{Correction nôtre, conforme à celle du critère (A) plus
haut. Contre-exemple : $\mathcal{C}_{n+1} = \mathrm{pt}$ et
$\mathcal{D}_{n+1} = K(A,n+1)$ ; $f_n$ est une équivalence entre points, et le
système local de fibre $K(A,n)$ donné par la fibration des chemins a des
sections sur $\mathcal{C}_{n+1}$ et aucune sur $\mathcal{D}_{n+1}$. Sur la
page, les accolades et les « $=$ déf » sous les deux membres sont ajoutés à
la main, et le second porte $\Gamma_{\mathcal{C}_{n+1}}(\mathcal{F})$ sans
$f_{n+1}^*$.}

\pagerange{54}{55}
\subsection*{App.~8 : trois approches de la cohomologie d'un topos (pages 54 et 55)}

La construction de la cohomologie par intégration de champs ne fait appel
ni aux complexes de faisceaux ni aux résolutions injectives ; elle
s'apparenterait plutôt aux calculs « čechistes » par hyperrecouvrements,
lesquels demandent un peu d'algèbre simpliciale. Il ne sait si une théorie
des champs pourra s'écrire sans algèbre simpliciale ; si oui, on aurait
trois approches :

\begin{itemize}
  \item[a)] complexes de faisceaux, résolutions injectives, catégories
  dérivées --- l'\emph{algèbre homologique commutative} : on résout les
  coefficients ;
  \item[b)] le point de vue čechiste ou simplicial --- l'\emph{algèbre
  homotopique} : on résout l'espace ;
  \item[c)] les $n$-champs --- l'algèbre catégorique ou \emph{algèbre
  homologique non commutative} : en apparence, on ne résout ni l'un ni
  l'autre.
\end{itemize}

La lettre se termine « Bien cordialement, Alexandre ».\footnote{Notre
commentaire. Dans un $\infty$-topos, c'est la définition c) qui est la
définition ordinaire : $H^n(X;A) = \pi_0\,\mathrm{Map}(1,K(A,n))$ (Lurie,
\emph{Higher Topos Theory}, § 7.2.2), et sa comparaison avec a) et b) est un
théorème --- pour b), le théorème des hyperrecouvrements de Verdier
(SGA~4, V.7) et sa version simpliciale (Dugger, Hollander, Isaksen, 2004). Mais
l'« apparence » qu'il note est juste : les fondements usuels des
$\infty$-catégories sont eux-mêmes simpliciaux, et la question de s'en
passer est celle que \emph{Pursuing Stacks} reprendra avec les
$\infty$-groupoïdes globulaires.}

\pagerange{56}{77}
\section*{Troisième lettre à Breen, Villecun, 17--19 juillet 1975 (pages 56 à 77)}

La lettre, qu'on lit dans la traduction anglaise revue en 1983 (voir la
page 37), répond à la réponse de Breen aux deux précédentes. Elle est datée
de sa main « Villecun 17/19 July 1975 ».\footnote{« Villecun » est lu avec
doute sur sa deuxième syllabe. Ses pages 2 à 11 manquent entre les pages 56
et 57 du dossier ; voir « Ce que le dossier n'établit pas ».}

\pagerange{56}{58}
\subsection*{App.~9 : une définition à trouver, et un Dold--Kan non abélien (pages 56 à 58)}

\emph{La définition des $n$-catégories.} Breen lui a proposé une construction
des $n$-catégories non strictes et du foncteur nerf. Il lui reconnaît « the
merit of existing » et d'être une première approche précise, mais la juge
très technique, peu intuitive alors qu'aux niveaux $1$ et $2$ « you just
follow your nose-tip », et surtout privée d'un foncteur des ensembles
simpliciaux vers les $n$-groupoïdes, sans lequel une théorie lui paraît
« kind of a joke ». D'où la méthode : écrire, « like a sort of list of
Christmas presents », toutes les notions qu'on voudrait avoir et les
relations qui doivent les lier, jusqu'à obtenir une description axiomatique
assez complète pour soit suggérer une construction \emph{ad hoc}, soit
démontrer un théorème d'existence et d'unicité d'une théorie du type
voulu.\footnote{La frappe porte « theorem », corrigé à la main en
« theory » ; l'appel (25) est ajouté à la main. Notre commentaire : c'est
la démarche des résultats d'unicité de la théorie des $(\infty,1)$-catégories
--- Toën, \emph{Vers une axiomatisation de la théorie des catégories
supérieures} (K-Theory, 2005), puis Barwick et Schommer-Pries,
\emph{On the unicity of the theory of higher categories} (J. Amer. Math.
Soc., 2021) ---, qui caractérisent la théorie par des axiomes, à l'automorphisme
« passage à l'opposée » près. La construction de Breen n'est pas dans ces
pages et on ne la décrit pas.}

\emph{Picard strict et non strict.} De la définition, qu'il n'a pas
comprise, des $n$-catégories de Picard proposée par Breen d'après Segal, il
ne retient qu'une exigence pour les \emph{strictes} : qu'elles forment une
$(n+1)$-catégorie $(n+1)$-équivalente à celle des complexes de chaînes de
longueur $n$. Et il remercie Breen d'avoir rectifié « a big blunder » : il
croyait que des H-espaces assez associatifs et commutatifs --- disons
équivalents à des $\Omega^iX$ pour $i$ arbitrairement grand ---
correspondaient à des groupes topologiques commutatifs.\footnote{L'erreur est
réelle et la correction de Breen est juste ; le commentaire qui suit est
nôtre. Un groupe topologique commutatif (ou un groupe abélien simplicial) a
le type d'homotopie d'un produit d'espaces d'Eilenberg--MacLane, alors
qu'un espace de lacets infini est un spectre connectif quelconque, dont les
invariants de Postnikov stables ne s'annulent pas en général. Les
$n$-catégories de Picard \emph{strictes} modélisent les complexes (les
$H\mathbb{Z}$-modules), les \emph{non strictes} les spectres connectifs
$n$-tronqués ; le premier écart est la multiplication par $\eta$, le signe
de la lettre à Deligne du dossier 134-1.}

\emph{Dold--Kan non abélien.} Reste une question : y a-t-il un analogue de
la théorie de Dold--Puppe pour les groupes simpliciaux non nécessairement
commutatifs ? Il tient le yoga
\[
\text{ensembles simpliciaux} \longleftrightarrow \infty\text{-groupoïdes}
\]
pour la version ultime « ensembliste » de Dold--Puppe, dont celle-ci se
déduirait en explicitant que les groupes abéliens dans les
$\infty$-catégories ne sont « rien d'autre » que les complexes de chaînes. Il
faut donc d'abord savoir ce que sont les \emph{groupes} dans les
$\infty$-catégories. Dans les $1$-catégories, il le sait --- ce sera
discuté, pense-t-il, dans la thèse de Mme Sinh, au chapitre des
Gr-catégories strictes ---, et donne la description en termes de théorie
des groupes : un quadruplet $(L_1,L_0,d,\theta)$, avec
$d : L_1 \to L_0$ un homomorphisme de groupes et
$\theta : L_0 \to \mathrm{Aut}(L_1)$ une action, soumis à\footnote{Dans (b), la frappe porte « $\theta(d(x_1) = \mathrm{int}(x_1)$ »,
sans la parenthèse fermante ; on l'a ajoutée.}
\[
\text{(a)}\quad d\bigl(\theta(x_0)\,x_1\bigr) = x_0\,d(x_1)\,x_0^{-1}, \qquad
\text{(b)}\quad \theta\bigl(d(x_1)\bigr) = \mathrm{int}(x_1).
\]
C'est ce qu'on appelle un \emph{module croisé} (J.H.C.~Whitehead, 1949) ; (b)
est l'identité de Peiffer. Il en tire que $\mathrm{Im}\,d$ est distingué, que
$\pi_1 = \mathrm{Ker}\,d$ est central dans $L_1$, que $L_0$ opère sur
$\pi_1$ à travers $\pi_0 = \mathrm{Coker}\,d$, et que l'invariant principal
est l'invariant de « Postnikov--Sinh » $\alpha \in H^3(\pi_0,\pi_1)$ ---
la classification des $2$-types connexes pointés par Mac Lane et Whitehead
(1950).\footnote{L'équivalence entre groupes dans les catégories (groupes
catégoriques stricts) et modules croisés est publiée par Brown et Spencer
(1976) ; la thèse de Hoàng Xuân Sính (1975) classe les Gr-catégories par le
même $H^3$. Le tapuscrit de cette thèse forme le dossier 135. Commentaire
nôtre.} Il a rencontré ces animaux à propos de la classification des groupes
formels « ordinaires » sur un corps parfait, à travers des groupes algébriques
affines et des groupes formels commutatifs reliés par des Gr-structures
strictes de ce type, dans un topos arbitraire, pour expliciter un principe
qu'il attribue à Dieudonné : « le caractère transcendant d'un groupe formel
est concentré essentiellement dans les groupes formels commutatifs ».

La question est la généralisation aux groupes dans les $n$-catégories :
il attend un « complexe de chaînes non commutatif »
\[
L_n \to L_{n-1} \to \cdots \to L_1 \to L_0 \to 1
\]
muni de structures supplémentaires du type de $\theta$, qui exprime
exactement le type d'homotopie tronqué en dimension $\leqslant n$ d'un groupe
topologique, ou, ce qui revient au même, le type d'homotopie en dimension
$\leqslant n+1$ d'un espace connexe pointé. « Have you candidates to
propose? »\footnote{Notre commentaire. Les réponses venues depuis :
pour $n = 2$, les $2$-modules croisés de Conduché (1984) ; pour tout $n$, les
cat$^n$-groupes de Loday (1982) et les $n$-cubes croisés d'Ellis et Steiner
(1987), qui modélisent les $(n+1)$-types connexes, mais ne sont pas des
complexes ; et, sous la forme même d'un complexe de groupes non commutatifs
muni d'opérations supplémentaires, les complexes hypercroisés de Carrasco et
Cegarra (1991), obtenus à partir du complexe de Moore d'un groupe
simplicial. Le théorème de Kan --- les groupes simpliciaux modélisent les
espaces connexes pointés --- répondait déjà à la question « ensembliste ».
Nous ne prétendons pas que ces travaux procèdent de la lettre.}

\pagerange{58}{60}
\subsection*{App.~10 : les six opérations et l'homologie (pages 58 à 60)}

Aux réflexions « naïves » de Breen sur la bidualité et l'homologie, il
répond par une série de réminiscences, numérotées (A) à (F).\footnote{Un
mot rayé de la page 58 est noirci et illisible. Des astérisques de sa main
surmontent « biduality » et « homology » ; leur renvoi n'est pas sur la
page.}

\emph{(A) Les six opérations.} Le formalisme de $Rf_!$ et $Rf^!$, combiné à
$Rf_*$, $Lf^*$, $\otimes^{\mathbb{L}}$ et $R\mathcal{H}om$ --- « les six
opérations » ---, contient implicitement la définition de l'homologie et
son identité essentielle avec la cohomologie. On l'a pour les faisceaux
quasi cohérents (séminaire Hartshorne, LNM 20), pour les espaces
localement compacts (Verdier, Bourbaki 300), pour la cohomologie étale à
coefficients premiers aux caractéristiques (SGA~5), pour les faisceaux
cohérents sur les espaces analytiques (Ramis, Ruget, Verdier). Il reste à le
développer dans le cadre cristallin, en caractéristique $0$ pour les modules
stratifiés à singularités « à la Deligne », peut-être avec des structures de
Hodge sur $\mathbb{C}$, enfin pour les motifs ; il est convaincu qu'il existe
« about anywhere \ldots\ wherever there is a formalism of a cohomological
nature ».\footnote{Notre commentaire. La prédiction s'est largement
réalisée : $\mathcal{D}$-modules holonomes réguliers et correspondance de
Riemann--Hilbert (Kashiwara, Mebkhout, vers 1980--1984), modules de Hodge
mixtes de Saito (1988--1990), six opérations motiviques (Ayoub, 2007 ;
Cisinski et Déglise), et, plus récemment, la notion même de « formalisme à six
foncteurs » comme objet d'étude (Liu et Zheng, Mann, Scholze). La forme
« developped » de la frappe est lue avec doute.}

Pour $f : X \to e = \mathrm{Spec}\,k$ séparé de type fini et un anneau de
coefficients $\Lambda$ de torsion première à la caractéristique, le
complexe $f^!\Lambda_e$ joue le rôle du complexe des chaînes singulières à
coefficients dans $\Lambda$, et $Rf_!\,f^!\Lambda_e$ celui de l'homologie
$H_*(X/e)$ vis-à-vis des coefficients sur $e$.\footnote{Précision nôtre :
$Rf_!f^!\Lambda$ est l'homologie ordinaire, duale de la cohomologie à
supports compacts ; $Rf_*f^!\Lambda$ serait l'homologie de Borel--Moore.
La « justification par un ou deux trucs » qu'il épargne à Breen est la
dualité globale.} Il ajoute : il n'y a pas à tronquer ; c'est lié à la
bidualité pour les coefficients constructibles, à valeurs dans un complexe
dualisant $K_X = f^!K_e$, qui est, pour $\Lambda = \mathbb{Z}/n\mathbb{Z}$
autodualisant, le « complexe des chaînes singulières » --- bidualité
démontrée en admettant la résolution des singularités, ou sans elle si l'on
en croit ce que Deligne lui a dit ; on peut traiter de même
$\mathbb{Z}_\ell$ (Jouanolou, thèse non publiée) ; et pour $f : X \to S$
lisse, ou plus généralement « cohomologiquement localement trivial », on a
$H_*(X/S) = Rf_!f^!\Lambda_S$.\footnote{La bidualité sans résolution est
démontrée par Deligne dans SGA~4$\tfrac12$, [Th. finitude]. Les travaux de
Jouanolou sur la cohomologie $\ell$-adique sont dans SGA~5. Commentaire
nôtre.}

\emph{(B) L'autodualité des jacobiennes.} Artin et Mazur ont étudié
l'autodualité de la jacobienne d'une courbe relative $X/S$ ; c'est acquis
sur un anneau de valuation discrète à fibre générique lisse, la fibre
spéciale pouvant être très sauvage. Il s'en est servi dans SGA~7 pour
démontrer, dans le cas des jacobiennes, sa conjecture sur l'accouplement
entre groupes de composantes des modèles de Néron de deux variétés
abéliennes duales.\footnote{Notre commentaire, avec réserve : la
conjecture est démontrée pour un corps résiduel fini (McCallum, 1986) et,
sauf erreur, pour un corps résiduel parfait (T.~Suzuki) ; elle est fausse en
général pour des corps résiduels non parfaits (Bertapelle et Bosch, 2000).}

Puis un souvenir : vers la fin des années cinquante, au moment où « the
grand cohomological stuff » sortait de l'ombre, le cours de Serre sur la
théorie de Rosenlicht et Lang des jacobiennes généralisées et du corps de
classes géométrique, et plus tard sa théorie « géométrique » du corps de
classes \emph{local} par les groupes proalgébriques, l'ont fait réfléchir à
des formulations cohomologiques de ces résultats, de nature géométrique,
dont les résultats arithmétiques sur un corps $k$ quelconque se
déduiraient par descente de $\bar k$. Il en esquissa des projets dans des
lettres à Serre, dans trois directions au moins, sans jamais réussir à
« mobiliser quelqu'un » pour les développer.

\pagerange{60}{61}
\subsection*{App.~11 : le complexe des jacobiennes généralisées (pages 60 et 61)}

\emph{(C)} Pour $X$ de type fini sur un corps $k$, construire un complexe
$J_{*X/k}$ de jacobiennes généralisées, de longueur $\dim X$, formé de
groupes proalgébriques commutatifs affines sauf $J_0$, dont la partie
abélienne est celle de la jacobienne généralisée usuelle $\mathrm{Alb}_{X/k}$.
Sa construction, inspirée du complexe résiduel, passe par des jacobiennes
généralisées (au sens cohomologique) des localisés
$\mathrm{Spec}\,\mathcal{O}_{X,x}$. Son $\underline{H}_0$ est la jacobienne
généralisée de $X$ : il existe un morphisme $X \to \underline{H}_0$
universel pour les morphismes de $X$ dans les groupes localement
proalgébriques commutatifs ; pour $X$ connexe, $\underline{H}_0$ est une
extension de $\mathbb{Z}$ par un groupe proalgébrique.\footnote{La frappe
dit « homeomorphism » ; il s'agit d'un morphisme de schémas. L'alinéa
« N.B. » qui le contient est encadré à la main et renvoyé par une flèche
après la phrase suivante, qu'on a suivie.} Le rôle cohomologique de ce
complexe est celui d'un complexe d'\emph{homologie} :
\[
(*) \qquad H^i(X,G_X) \simeq \mathbb{E}\mathrm{xt}^i\bigl(J_{*X/k},G\bigr),
\]
pour des coefficients $G$ qu'il prenait alors parmi les groupes algébriques
commutatifs, mais avec la topologie de Zariski (« malédiction ! »), ce qui
donne des groupes « légèrement stupides ». Il faudrait reprendre la
construction en cohomologie fppf, avec une formule $(*)$ sans limitation sur
$i$ --- le complexe $J_*$ s'arrêtant alors en dimension $2\dim X$, ses
composantes au-delà de $\dim X$ pouvant être de dimension nulle --- et dans
une catégorie dérivée de groupes proalgébriques, notion que personne n'avait
définie quand Serre développait le formalisme de ces groupes. La
construction de $J_*$ ne commuterait au changement de base qu'au sens
dérivé.\footnote{L'indice de « $(J_\cdot)_{X/k}$ » est surchargé et lu avec
doute ; un mot biffé après « and » est illisible ; le numéro de l'intertitre
App.~11 est surchargé. Notre commentaire : l'objet le plus proche qu'on
connaisse aujourd'hui est le complexe d'Albanese motivique $\mathrm{LAlb}(X)$
de Barbieri-Viale et Kahn (\emph{On the derived category of 1-motives},
Astérisque 381, 2016), complexe de $1$-motifs défini pour tout $X$, avec les
jacobiennes généralisées « à module » de Kato et Russell ; nous ne disons
pas qu'ils réalisent le programme de la lettre.}

\pagerange{62}{65}
\subsection*{App.~12 : le corps de classes géométrique global comme dualité (pages 62 à 65)}

\emph{(D)} Soit $f : X \to \mathrm{Spec}\,k$ lisse, séparé, de type fini, de
dimension relative $d$, et $n > 0$ premier à la caractéristique. Pour $F$
annulé par $n$, la dualité globale dit que $Rf_!F$ et
$Rf_*R\mathcal{H}om(F,\mu_n^{\otimes d})$ sont duaux à valeurs dans
$(\mathbb{Z}/n\mathbb{Z})_k$, avec un décalage de $2d$ ; comme
$\mathbb{Z}/n\mathbb{Z}$ est injectif sur lui-même, cela donne des dualités
parfaites\footnote{La frappe donne pour second exemple « $Rf_!(\mu_m^{\otimes})$ et
$Rf_*(\mathbb{Z}/n\mathbb{Z})$ » ; il faut lire
$Rf_!(\mu_n^{\otimes d})$.}
\[
R^if_!(F) \times R^{2d-i}f_*\bigl(R\mathcal{H}om(F,\mu_n^{\otimes d})\bigr) \longrightarrow \mathbb{Z}/n\mathbb{Z}.
\]
Si $n$ est une puissance de $p > 0$, tout
s'effondre : on ne sait même plus, pour $d > 1$, par quoi remplacer
$\mu_n^{\otimes d}$.\footnote{Notre commentaire : la réponse est venue
quelques années plus tard, avec les faisceaux de de Rham--Witt logarithmiques
$W_r\Omega^d_{X,\log}$, placés en degré $d$ (Milne, Gros, Kato, puis Moser
pour les coefficients constructibles), qui jouent le rôle de
« $\mu_{p^r}^{\otimes d}$ ».}

« The extraordinary miracle » est que pour $d = 1$, une courbe lisse, tout
continue de marcher, pourvu qu'on énonce avec soin : on le vérifie sur
$\mathbb{Z}/p\mathbb{Z}$, $\mu_p$, $\alpha_p$ pour $X$ complète, par
l'autodualité de la jacobienne, puis pour des coefficients tordus et $X$ non
complète. Seulement, les $R^if_*F$ et $R^if_!F$ ont alors une structure
« continue » : ce sont essentiellement des groupes proalgébriques, comme le
$\pi_1^{\mathrm{ab}}$ d'une courbe non complète, donc
$H^1(X,\mathbb{Z}/p^n\mathbb{Z})$.\footnote{Dans
« $H^1(X,\mathbb{Z}/p^n\mathbb{Z})$ » la frappe portait
$\mathbb{Z}/n\mathbb{Z}$, corrigé à la main, et corrigé de nouveau en marge.}

\emph{La dualité de Serre.} Serre avait proposé, sans l'écrire, une dualité
pour les groupes algébriques unipotents commutatifs \emph{à isogénie
radicielle près} (« quasi-algébriques ») : sur $k$ algébriquement clos,
$G' = \mathrm{Ext}^1(G,\mathbb{Q}/\mathbb{Z})$ porte une structure canonique
de groupe quasi-algébrique, avec
$\mathrm{Ext}^1(\mathbb{G}_a,\mathbb{Q}/\mathbb{Z}) \simeq \mathbb{G}_a$, et
$\Delta G = G'$ vérifie $G \simeq \Delta\Delta G$ --- « an authentic
autoduality », qu'il appelle \emph{dualité de Serre}. Pour les groupes
finis étales, c'est $\mathrm{Ext}^0(G,\mathbb{Q}/\mathbb{Z})$, la dualité de
Pontrjagin, qui est parfaite. On les recolle dans une catégorie dérivée
convenable en posant
\[
\Delta G = R\mathcal{H}om(G,\mathbb{Q}/\mathbb{Z}),
\]
dualité « cohomologique » de Serre, autodualité parfaite sur les complexes
dont les $\underline{H}^i$ sont, à passage à la limite près, extensions de
groupes étales par des groupes unipotents connexes --- les seuls qu'on
rencontre en intégrant des coefficients finis par $Rf_!$ ou
$Rf_*$.\footnote{Notre commentaire. Les groupes quasi-algébriques de Serre
sont aujourd'hui les schémas en groupes \emph{parfaits}, et cette dualité
est établie dans la catégorie dérivée des faisceaux sur le site parfait
(Bégueri, 1980 ; Milne, \emph{Arithmetic Duality Theorems}, chapitre III) ;
le théorème d'annulation de Breen pour les $\mathrm{Ext}$ de faisceaux
abéliens (1975) en est un ingrédient. Sur la page, l'exposant frappé de
$G' = \mathrm{Ext}^1$ est biffé et remplacé par un accent, répété en
marge.} Passant à la limite sur $n$, avec $f^!(\mathbb{Q}/\mathbb{Z})_k
= (\mu_\infty)_X[2]$, la formule de dualité pour une courbe devient
\[
\Delta\bigl(Rf_!F\bigr) \simeq Rf_*\bigl(DF[2]\bigr),
\]
où $D = R\mathcal{H}om(-,\mu_\infty)$ est la dualité de Cartier (ou
$R\mathcal{H}om(-,\mathbb{G}_m)$) et $\Delta$ la dualité de Serre :
cohomologie à supports propres et cohomologie ordinaire s'échangent, « when
one takes upstairs Cartier duality, and downstairs Serre
duality ».\footnote{La page écrit
« $f^!(\mathbb{Q}/\mathbb{Z})_k = (\mu_\infty)_X$ » sans décalage ; le
décalage $[2]$ est celui de la formule, et c'est à lui que renvoie, sur la
page, une flèche tracée depuis le mot « shift ». L'indice $\infty$ de
$\mu_\infty$ est écrit à la main sous un indice frappé noirci.}

La validité de la formule « is not open to doubt » ; le travail est de
décrire soigneusement les catégories de coefficients, sur $X$ et sur $k$, et
les foncteurs $D$ et $\Delta$ ; la flèche étant immédiate, la formule
résulterait des dévissages usuels vers $\mathbb{Z}/p\mathbb{Z}$, $\mu_p$,
$\alpha_p$ sur $X$ complète. Appliquée à $R^1f_*(G_X)$ pour $G$ fini étale
sur $k$, elle redonne la description de Serre du corps de classes
géométrique en termes d'extensions par $G$ d'une jacobienne généralisée :
\emph{la formule de dualité est une version cohomologique, considérablement
enrichie, du corps de classes géométrique.}\footnote{Dans « $G_X$ », la
frappe porte le même $G$ évidé que dans $\mathbb{G}_a$ ; d'après la suite, il
s'agit du faisceau constant de valeur $G$.} Sur un corps fini, on retrouve
la forme classique par « le truc de Lang » --- l'isogénie de Lang identifie
le $\pi_1$ arithmétique d'un groupe algébrique commutatif lisse connexe $J$
à $J(k)$ ---, sous la forme cohomologique
\[
\Delta_0\,R\Gamma_k(J^*) \simeq R\Gamma_k\bigl(\Delta J^*[1]\bigr),
\]
où $\Delta_0$ est la dualité de Pontrjagin des groupes topologiques totalement
discontinus et $J^*$ un complexe d'ind-progroupes algébriques ; en
l'appliquant à la formule géométrique, on obtient la formule de dualité du
corps de classes arithmétique\footnote{Notre commentaire. C'est, pour une courbe sur un corps fini et des
schémas en groupes finis plats, la dualité d'Artin--Milne
(\emph{Duality in the flat cohomology of curves}, Invent. Math., 1976 ;
Milne, \emph{Arithmetic Duality Theorems}, III) :
$H^i_c(X,F)\times H^{3-i}(X,F^D)\to\mathbb{Q}/\mathbb{Z}$ parfait. La
version « géométrique », avec la dualité de Serre en bas, est le programme
que T.~Suzuki a repris par des sites relatifs ; nous ne disons pas
jusqu'où.},
\[
\Delta_0\bigl(H^*_!(X,F)\bigr) \simeq H^*\bigl(X,D(F)[3]\bigr).
\]

Enfin, pour un $F$ qui n'est pas étale mais a une structure continue, comme
$\alpha_p$, il faut définir $Rf_!F$ avec soin à partir d'une
compactification $\tilde f : \tilde X \to \mathrm{Spec}\,k$, au moyen d'un
triangle exact

\begin{tikzcd}[column sep=small]
 & R\hat{f}_*\bigl(\hat{\tilde F}\bigr) \arrow[dl, "{+1}"'] & \\
Rf_!(F) \arrow[rr] & & R\tilde{f}_*(\tilde{F}) \arrow[ul]
\end{tikzcd}

où $\hat{\tilde F}$ est le complété formel d'un prolongement
« admissible » $\tilde F$ de $F$ le long des points à
l'infini.\footnote{La page note le sommet « $Rf_!(\hat{\tilde F})$ » ; il
s'agit de l'image directe, sur le spectre de $k$, de la cohomologie du
voisinage formel de $\tilde X \setminus X$, ce qu'on a écrit
$R\hat f_*$. Le degré de la flèche de bord, implicite sur la page, est
marqué $+1$. Notre commentaire : c'est ainsi que Milne définit la
cohomologie à supports compacts en cohomologie plate, par les complétés aux
places manquantes (\emph{Arithmetic Duality Theorems}, III, § 0).} C'est là,
dit-il, qu'apparaît le lien avec le corps de classes local.

\pagerange{65}{66}
\subsection*{App.~13 : le corps de classes local comme dualité (pages 65 et 66)}

\emph{(E)} Soit $V$ un anneau de valuation discrète complet, de corps
résiduel $k$ et de corps des fractions $K$ ; on suppose soit $k$ relevé dans
$V$ (donc $V \simeq k[[T]]$), soit $k$ parfait de caractéristique $p>0$.
Pour des coefficients $F$ finis sur $K$, le travail principal serait de
définir une catégorie convenable de coefficients sur $k$ et un foncteur
$F \mapsto R\underline{\Gamma}_K(F)$ à valeurs dans celle-ci, avec
\[
R\Gamma_K(F) \simeq R\Gamma_k\bigl(R\underline{\Gamma}_K(F)\bigr),
\]
conformément à l'intuition que, pour $k$ algébriquement clos, les
$H^i(K,F)$ portent une structure de groupe algébrique (ind-pro) sur $k$ ;
les vecteurs de Witt et le foncteur de Greenberg y joueraient un rôle
essentiel. La formule de dualité s'énoncerait alors comme en (D),
\[
\Delta\,R\underline{\Gamma}_K(F) \simeq R\underline{\Gamma}_K\bigl(DF[1]\bigr),
\]
et deviendrait, pour un corps résiduel fini, par le truc de Lang,
\[
\Delta_0\,R\Gamma_K(F) \simeq R\Gamma_K\bigl(DF[2]\bigr),
\]
qui contient le corps de classes local géométrique de Serre et le corps de
classes local arithmétique classique.\footnote{La frappe porte $\Gamma_k$ au
premier membre de cette dernière formule et $\Gamma_K$ au second ; c'est
$\Gamma_K$ des deux côtés, et c'est alors la dualité locale de Tate,
$H^i(K,F)\times H^{2-i}(K,F^D)\to\mathbb{Q}/\mathbb{Z}$. L'exposant de
$H^i(K,F)$ est surchargé et lu avec doute.}

Remarques. (a) Pour $F$ premier à la caractéristique résiduelle, c'est facile
et connu : un cas particulier de la formule d'induction
$i^!\,D_S(F) \simeq D_s(i^*F)$ pour l'immersion fermée
$i : s = \mathrm{Spec}\,k \to S = \mathrm{Spec}\,V$ ; le travail porte donc
sur les coefficients $p$-primaires, et le cas le plus subtil est
l'inégale caractéristique.\footnote{La page écrit « $D_S(i^*(F))$ » au second
membre ; c'est le foncteur dualisant de $s$ qui convient, $i^*F$ vivant sur
$s$. Elle écrit aussi « $p = \mathrm{Spec}(k)$ » pour le point fermé.}
(b) $R\underline{\Gamma}_K$ s'obtiendrait en composant $Rj_*$, pour
$j : \mathrm{Spec}\,K \to \mathrm{Spec}\,V$, avec une version cohomologique
du foncteur de Greenberg. (c) Il ne parle que de coefficients finis, les
seuls pour lesquels il est sûr de ce qu'il avance ; mais il se rappelle
vaguement, sans garantie, une formule pour un schéma abélien $F$ sur $K$, de
dual $F'$, avec $G'$ le groupe proalgébrique attaché « à la Greenberg » au
modèle de Néron de $F'$ ($k$ algébriquement clos) :
\[
H^1(K,F) \overset{?}{\simeq} \mathrm{Ext}^1_{k\text{-grp}}(G',\mathbb{Q}/\mathbb{Z}),
\]
et pense que la conjecture de SGA~7 sur les modèles de Néron devrait sortir
de cette machine de dualité locale.\footnote{Notre commentaire, avec
réserve. Pour $G'$ connexe, $\mathrm{Ext}^1(G',\mathbb{Q}/\mathbb{Z})$ est le
dual de $\pi_1(G')$, et la formule est la conjecture de Šafarevič, démontrée
par Bégueri en caractéristique nulle et par Bertapelle en égale
caractéristique $p$. Le numéro « SGA~7 » est surchargé à la main sur un
chiffre frappé, peut-être 6.} (d) « You may ask Deligne if he didn't dive
into questions (D) and (E) lately. »

\pagerange{67}{67}
\subsection*{App.~14 : portée et limites de la topologie fppf (page 67)}

\emph{(F)} Depuis les tentatives de Serre pour trouver une cohomologie de
Weil avec des coefficients continus (les vecteurs de Witt $W_n$), il a eu
l'impression qu'une cohomologie de Weil $p$-adique correcte en
caractéristique $p$ devrait venir de la cohomologie fppf, à coefficients
finis ou dans des groupes algébriques sur $k$ ; la construction du complexe
jacobien local était liée à cet espoir --- « the homology might reveal what
is hidden to us in cohomology! ».\footnote{La page dit « The construction in
(B) of the local jacobian complex » ; c'est en (C) que la lettre construit
ce complexe, et (B) porte sur l'autodualité des jacobiennes.} Mais on a
maintenant la cohomologie cristalline, qui a les bonnes propriétés pour $X$
projective et lisse (Berthelot), et, prise pour étalon, elle montre que la
partie de $H^i_{\mathrm{cris}}(X)$ descriptible par la cohomologie fppf à
coefficients dans des $k$-groupes algébriques n'en est qu'une petite part :
celle des pentes $\leqslant 1$ du Frobenius, alors que les pentes vont de $0$
à $i$. D'où un $H^1$ correct, et l'échec pour $H^2$. On conjecture que
\emph{toute} la partie de pentes $\leqslant 1$ vient de fppf ; Mazur, Katz,
Messing et Deligne en sauraient l'état.\footnote{Notre commentaire. La
décomposition par les pentes a été rendue précise par le complexe de
de Rham--Witt (Illusie, 1979) : la partie de pentes $< 1$ est
$H^i(X,W\mathcal{O}_X)_{\mathbb{Q}}$, la cohomologie de Witt de Serre, et la
partie de pente $1$ est atteinte par la cohomologie plate de $\mu_{p^\infty}$,
ou de $W_r\Omega^1_{\log}$. Les coefficients de Tate supérieurs
$\mathbb{Z}_p(n)$ de la cohomologie syntomique, puis prismatique
(Bhatt, Morrow, Scholze ; Bhatt, Lurie), atteignent les pentes jusqu'à $n$ ;
c'est une manière de dépasser la limite qu'il constate, sans que nous
prétendions qu'elle réalise son espoir. L'indice de « $\mu_{p^n}$ » est
surchargé et corrigé en marge.}

\pagerange{67}{69}
\subsection*{App.~15 : le type d'homotopie d'un topos (pages 67 à 69)}

La question 7 de Breen trahit, dit-il, un malentendu sur le sens du « type
d'homotopie » d'un topos $X$ : on confond l'algèbre homotopique qu'on peut
faire \emph{sur} $X$ --- faisceaux simpliciaux, champs de toutes sortes ---
et le point de vue où $X$, avec sa riche structure, disparaît pour ne plus
être qu'un « pale element » d'une catégorie homotopique (ou
pro-homotopique), obtenu par une localisation très poussée. Il ne reste
alors à $X$ que ses $\pi_i$ et sa cohomologie à coefficients constants ou
tordus ; et, à creuser ce qui reste, on tombe sur \emph{les $n$-champs
localement constants} --- une équivalence d'homotopie $f : X \to X'$
induisant une $(n+1)$-équivalence entre les $(n+1)$-catégories des
$n$-champs localement constants --- qui contiennent les complexes abéliens
de longueur $n$ à faisceaux de cohomologie localement constants. D'où un
triangle d'objets qui se déterminent mutuellement, les deux sommets du bas
étant pris à $n$-équivalence près, et le sommet du haut pouvant être aussi
un ensemble simplicial :

\begin{tikzcd}[column sep=small, row sep=large, nodes={font=\scriptsize}]
 & \text{topos (ou espace) à $n$-homotopie près} \arrow[dl, leftrightarrow] \arrow[dr, leftrightarrow] & \\
\text{$n$-groupoïdes} \arrow[rr, leftrightarrow] & & \text{$n$-catégories « spéciales »}
\end{tikzcd}

Une $n$-catégorie $E_n$ est « spéciale », ou \emph{$n$-galoisienne},
si elle est $n$-équivalente à celle des $(n-1)$-champs localement constants
sur un espace ou un topos, ou, ce qui doit revenir au même, à celle des
$n$-foncteurs $G_n \to ((n-1)\text{-}\mathrm{Cat})$ pour un $n$-groupoïde
$G_n$ ; $G_n$ est alors le $n$-groupoïde fondamental, $E_n$ la catégorie des
$(n-1)$-systèmes locaux, $X$ la réalisation géométrique. Comme pour $n=1$,
$G_n$ devrait se retrouver comme la sous-$n$-catégorie pleine de
$\underline{\mathrm{Hom}}_n(E_n,((n-1)\text{-}\mathrm{Cat}))$ formée des
$n$-foncteurs ayant des propriétés d'exactitude --- « which commute with
finite $\varprojlim$ and arbitrary $\varinjlim$ » ---, ce qui soulève « the
disquieting vision of $n$-limits in $n$-categories ».\footnote{La page écrit
« $n$-functors $G_n \to (n\text{-}\mathrm{Cat})$ » et
« $\mathrm{Hom}_n(E_n,(n\text{-}\mathrm{cat}))$ », alors que, d'après sa
propre correction de « local $n$-systems » en « local $(n-1)$-systems », il
s'agit de $(n-1)$-champs : avec $(n\text{-}\mathrm{Cat})$ pour but, $E_n$
serait une $(n+1)$-catégorie. On a mis $((n-1)\text{-}\mathrm{Cat})$, comme
dans la lettre du 17 février. Le triangle occupe sur la page plusieurs lignes
par sommet ; on l'a traduit. Notre commentaire : c'est, à la lettre, la
théorie de Galois supérieure de Hoyois citée page 49 ; les $n$-limites dans
les $n$-catégories sont devenues ordinaires avec les $\infty$-catégories.}

Il est prudent de supposer $X$ localement homotopiquement trivial, ce qui
assure que le pro-ensemble simplicial d'Artin--Mazur est essentiellement
constant ; ce n'est sûrement pas le cas du topos étale d'un schéma, dont le
$n$-groupoïde fondamental doit être un pro-$n$-groupoïde, et $E_n$ une
ind-$n$-catégorie, l'ind-structure traduisant la trivialité locale relative
à des recouvrements de plus en plus fins.

\pagerange{69}{72}
\subsection*{App.~16 : stratifications et voisinages tubulaires (pages 69 à 72)}

Il comprend pourtant la résistance de Breen à ce dépouillement d'un « beau
topos », et pense qu'en allant à sa racine on arrive à approfondir la notion
même de type d'homotopie.

\emph{Faisceaux $P$-constructibles.} Sur le topos étale d'un schéma
cohérent $X$, les faisceaux constructibles sont ceux pour lesquels il existe
une partition finie $P$ de $X$ en strates localement fermées constructibles
$X_i$ telle que la restriction à chaque $X_i$ soit localement constante.
Pour $P$ fixée, les faisceaux (ou complexes, ou champs) $P$-constructibles
ne forment une catégorie satisfaisante que si elle est stable par les
opérations usuelles, $R\mathcal{H}om$, $Rj_*j^*$ pour une strate $j$, etc. ;
pour $X$ excellent et avec la résolution des singularités, les faisceaux
constructibles de torsion première aux caractéristiques sont stables, mais
pas pour une partition fixée. Il faut pour cela des hypothèses très strictes
d'« équisingularité » le long des strates, et il pense qu'un raffinement des
techniques connues fournira des stratifications équisingulières
arbitrairement fines.

\emph{Deux strates.} Soit $X_0$ fermé, $X_1 = X \setminus X_0$, avec les
inclusions $X_0 \xrightarrow{i_0} X \xleftarrow{i_1} X_1$. Par le dévissage
d'Artin (le recollement de SGA~4), un faisceau $F$ sur $X$ équivaut à la
donnée de $F_0 = i_0^*F$, de $F_1 = i_1^*F$ et d'un morphisme
$F_0 \to \varphi(F_1)$, avec $\varphi = i_0^*\,i_{1*}$ exact à gauche. $F$
est $P$-constructible si et seulement si $F_0$ et $F_1$ sont localement
constants ; par équisingularité, $\varphi(F_1)$ l'est aussi, et tout
s'exprime par les faisceaux localement constants sur $X_0$ et $X_1$ --- leur
seul type d'homotopie --- et le foncteur $\varphi$. Pour expliciter
celui-ci, il propose un « voisinage tubulaire étale » de $X_0$, topos
intéressant mais non associé à un schéma, paraphrase technique d'une intuition
topologique très simple.\footnote{La page dit « tubular neighbourhood of
$X_0$ in $X_1$ » ; il s'agit d'un voisinage de $X_0$ dans $X$, dont on
retient la trace sur $X_1$. Un mot biffé avant « $\varphi(F_1)$ » est
noirci ; « topological intuition » est entouré et renvoyé après « simple ».}

\emph{L'intuition.} Sur $\mathbb{C}$, l'équisingularité dit qu'il existe un
voisinage tubulaire $T$ de $X_0$ dans $X$, rétracté sur $X_0$, tel que
$T\setminus X_0$ soit localement trivial sur $X_0$ ; si $\partial T$ est son
bord, $T$ est le fibré en cônes
$(\partial T \times [0,1]) \amalg_{\partial T} X_0$, et
$\mathring T = T \setminus X_0 \simeq \partial T \times [0,1[$ est
$X_0$-homotope à $\partial T$. Si les strates sont non singulières,
$\partial T \to X_0$ est une fibration topologiquement lisse ; avec
$\widetilde X_1 = X_1 \setminus \mathring T$, l'inclusion
$\widetilde X_1 \to X_1$ est une équivalence d'homotopie, et $X$ se
reconstitue à homéomorphisme près comme somme amalgamée du diagramme

\begin{tikzcd}
\partial T \arrow[r, "j"] \arrow[d, "p"'] & \widetilde{X}_1 \\
X_0 &
\end{tikzcd}

où $p$ est une fibration et $j$ l'inclusion du bord. Alors
$\varphi(F_1) \simeq p_*j^*\widetilde F_1$, et la donnée de $F$ est celle
de
\[
F_0, \qquad \widetilde F_1, \qquad \tilde u : p^*F_0 \to j^*\widetilde F_1,
\]
avec $F_0$, $\widetilde F_1$ localement constants sur $X_0$ et
$\widetilde X_1$.\footnote{La page note $\phi$ et $\varphi$ le même foncteur ;
on a gardé $\varphi$. Sur la page, $\mathring T$ est « $X_0$-homotopic » à
$\partial T$, et le crochet de « $[0,1[$ » est corrigé à la main.} Il
propose aussi les diagrammes équivalents
$X_0 \leftarrow \mathring T \to X_1$, plus canonique, et
$T \hookleftarrow \mathring T \to X_1$, où l'inclusion
$\mathring T \to T$ est « morally » une fibration à fibres compactes
non singulières de dimension finie --- bien mieux que ce que donne le yoga
de Cartan--Serre selon lequel toute application est équivalente à une
fibration ---, et qui a l'avantage de se construire algébriquement pour les
schémas, une fois les voisinages tubulaires étales
construits.\footnote{Notre commentaire. Ce dictionnaire est aujourd'hui la
théorie des \emph{chemins de sortie} : pour un espace conique stratifié, les
faisceaux constructibles (d'ensembles, de complexes ou d'espaces) sont les
foncteurs sur l'$\infty$-catégorie des chemins de sortie (MacPherson ;
Treumann, \emph{Exit paths and constructible stacks}, 2009 ; Lurie,
\emph{Higher Algebra}, A.9), et, pour deux strates, ce foncteur est
exactement la donnée $(F_0, F_1, p^*F_0 \to j^*F_1)$, le « lien » $\partial T$
portant les chemins qui sortent de $X_0$. La version étale, pour les
schémas, est l'« exodromie » de Barwick, Glasman et Haine (2018). Les
voisinages tubulaires algébriques ont été étudiés par Cox (1979).}

\pagerange{72}{75}
\subsection*{App.~17 : le type d'homotopie fin (pages 72 à 75)}

Le point est que les faisceaux, complexes ou $n$-champs $P$-constructibles,
pour une stratification équisingulière $P$, se ramènent à la connaissance
d'un diagramme de \emph{types d'homotopie} ordinaires (des pro-types en
topologie étale), en prenant des systèmes locaux sur les sommets reliés par
des morphismes de compatibilité du type $p^*F_0 \to j^*F_1$.\footnote{Le
sommet du diagramme de la page 73 est frappé « $\Delta$ », pour
$\partial T$ ; le numéro de l'intertitre App.~17 est surchargé et lu avec
doute.} Ce serait un exercice amusant, qu'il n'a pas fait, de traduire les
six opérations dans ce dictionnaire et d'y retrouver les formules connues,
au moins pour des strates non singulières ; mais il faudrait sans doute
considérer non un diagramme dans la catégorie homotopique, mais la catégorie
homotopique des diagrammes d'ensembles simpliciaux.\footnote{C'est
exactement le point technique qu'on a retenu depuis : un diagramme dans la
catégorie homotopique perd l'information de cohérence, et c'est
l'$\infty$-catégorie des chemins de sortie, ou la catégorie de modèles des
diagrammes, qu'il faut prendre. Commentaire nôtre.}

\emph{La topologie modérée.} En réfléchissant aux fondements d'une « tame
topology », qui élimine d'emblée les phénomènes sauvages, il a récemment
explicité comment une stratification équisingulière à strates non
singulières d'un espace modéré compact donne canoniquement un diagramme
fini de variétés à bord, reliées par des fibrations localement triviales de
variétés à bord compactes et par les inclusions des bords --- avec des bords
munis d'une décomposition cellulaire élémentaire ---, à partir duquel
l'espace se reconstitue par recollement. Pour un morphisme modéré
$f : X \to Y$ et des stratifications adaptées, on obtiendrait un morphisme
de diagrammes, et les quatre opérations $Rf_*$, $Rf_!$, $Lg^*$, $g^!$
entre faisceaux $P_X$- et $P_Y$-constructibles s'exprimeraient par des
opérations standard entre types d'homotopie. Le tout devrait se transposer
aux schémas excellents par les voisinages tubulaires étales.\footnote{Notre
commentaire. La « topologie modérée » a reçu une forme dans les structures
o-minimales (van den Dries, \emph{Tame topology and o-minimal structures},
1998, dont le titre reprend l'expression de l'\emph{Esquisse d'un
programme}) ; la décomposition en variétés à bord à partir d'une
stratification équisingulière est celle de la théorie de Thom--Mather
(données de contrôle, lemmes d'isotopie).}

\emph{Le type d'homotopie fin} d'un espace modéré ou d'un schéma excellent se
définit alors par passage à la limite sur les « $P$-types d'homotopie »
associés à des stratifications équisingulières de plus en plus fines. Il
contiendrait la connaissance non seulement des faisceaux localement
constants, mais, à la limite inductive, de \emph{tous}, et dépendrait
fonctoriellement de $X$. Pour un schéma de type fini sur un corps
algébriquement clos, le théorème de finitude le plus fort s'exprimerait par
ce type fin : ses constituants seraient des « polyèdres finis », voire des
variétés compactes à bord, ou plus précisément leurs complétés profinis
premiers à $p$ seraient ceux de tels polyèdres. On voit comment commencer en
caractéristique $0$ ; en caractéristique $p$, il prévoit un supplément
d'amusement, voire de mystère, pour les variétés qui résistent à tout
relèvement en caractéristique $0$, même birationnel.\footnote{Notre
commentaire. Le type d'homotopie fin est, sous un autre nom, la « catégorie
de Galois » profinie stratifiée de l'exodromie de Barwick, Glasman et Haine,
qui code tous les faisceaux constructibles d'un schéma cohérent, et la
théorie de l'homotopie des espaces stratifiés qui l'accompagne (Haine). Nous
ne savons pas dire où en est l'énoncé de finitude en caractéristique $p$, et
le laissons comme la lettre le laisse. Sur la page, « the ordinary homotopy
types … finite polyhedra » est souligné à la main.}

\pagerange{75}{77}
\subsection*{App.~18 : linéariser un type d'homotopie (pages 75 à 77)}

De ces pensées géométriques il ne tire pas de programme précis, et se borne
à des remarques. Il se demande depuis longtemps combien d'information
cohomologique il faut pour reconstituer un type d'homotopie : le groupe
fondamental, les foncteurs $H^i(-,M)$ sur les $\pi_1$-modules avec leur
structure de foncteurs cohomologiques, les cup-produits, d'autres
opérations. Avec le langage des catégories dérivées, un meilleur candidat
est la sous-catégorie de la catégorie dérivée des complexes abéliens sur $X$
à faisceaux de cohomologie localement constants, avec sa structure
triangulée et son produit $\otimes^{\mathbb{L}}$, peut-être
$R\mathcal{H}om$ ; il ne sait pas si cela suffit. En revanche il ne doute pas
qu'en poussant la linéarisation jusqu'au bout, c'est-à-dire jusqu'au cadre
\emph{non abélien} --- les $(n+1)$-catégories des $n$-champs localement
constants, sans aucune structure supplémentaire ---, on reconstitue le type
d'homotopie par son $\infty$-groupoïde fondamental, ce qui signifie que
toutes les opérations cohomologiques imaginables sont déjà dans ces
données.\footnote{Notre commentaire. La version non abélienne est
aujourd'hui une tautologie bien comprise : l'$\infty$-topos
$\mathrm{Fun}(X,\mathcal{S})$ des systèmes locaux d'espaces détermine $X$
comme sa forme. Pour la version abélienne, la réponse positive la plus
nette porte sur les cochaînes munies de leur structure d'algèbre
$E_\infty$ : Mandell (\emph{E$_\infty$ algebras and p-adic homotopy theory},
2001 ; \emph{Cochains and homotopy type}, 2006) montre qu'elles déterminent
les espaces nilpotents de type fini, $p$-complétés ou entiers ; la
structure triangulée et $\otimes^{\mathbb{L}}$ seuls ne suffisent pas en
général, faute de cohérences. L'appel (31) est répété en marge.}

De même, les types d'homotopie plus élaborés attachés aux stratifications
devraient correspondre aux $(n+1)$-catégories des $n$-champs subordonnés à
$P$ ; et, si la description par la catégorie dérivée « localement
constante » valait, on retrouverait le type d'homotopie mixte par la
sous-catégorie dérivée des complexes à cohomologie localement constante sur
chaque strate, avec $\otimes^{\mathbb{L}}$, $R\mathcal{H}om$ et les quatre
opérations entre réunions localement fermées de strates. Le problème est
qu'on ne sait même pas ce qu'est une catégorie triangulée, ni sa version
non commutative, décrite sans doute plus simplement et plus
fondamentalement comme une « catégorie homotopique » munie des opérations
de prise de fibres et de cofibres.\footnote{Une parenthèse biffée de
traits obliques, page 77, ajoutait : « (N.B. A first approach has been
developed by Quillen at the beginning of the years “70”.) » ;
\emph{Homotopical Algebra} (LNM 43) est de 1967. Notre commentaire : la
« catégorie homotopique avec fibres et cofibres » est aujourd'hui une
$\infty$-catégorie à limites et colimites finies, dont la version stable
est la « bonne » catégorie triangulée ; c'est aussi ce que les dérivateurs
de \emph{Pursuing Stacks} cherchent à capter.}

« It is surely time that I finish this “lettre-fleuve”, which is becoming
more and more vague. » Une dernière question : quelle est cette « marvellous
formula of Bloch--Quillen » à laquelle Breen fait allusion, dont il n'a
jamais entendu parler, « and which makes my mouth water? »\footnote{C'est,
selon toute vraisemblance, la formule de Bloch
$\mathrm{CH}^p(X) \simeq H^p(X,\mathcal{K}_p)$, démontrée par Bloch pour
$p = 2$ (1974) et par Quillen en général pour les variétés lisses (1973).
L'allusion de Breen n'est pas dans ces pages ; l'identification est la
nôtre. La pagination de la dernière page, surchargée, est lue avec doute
« 61 bis ».} Signé : « Very cordially yours, Alexandre Grothendieck. »

\pagerange{78}{80}
\section*{Les notes au chapitre I et à l'appendice (pages 78 à 80)}

Les pages 78 à 80 sont les trois premières pages du tapuscrit « NOTES to
Chapter I and Appendix », notes (1) à (25), avec des corrections et ajouts
de sa main, dont la note (11), « Added 23.2.83 ». Elles commentent les
sections du chapitre I et, à partir de la note (22), les lettres à Breen.
Ce texte appartient à \emph{Pursuing Stacks} et il est imprimé dans
G.~Maltsiniotis (éd.), \emph{À la poursuite des champs}, vol.~I, SMF,
Documents mathématiques 20, 2022 ; on y renvoie et on ne le restitue pas. La
note (25) se poursuit au-delà de la page 80, dans la partie du dossier qui
n'est pas transcrite.\footnote{La transcription relève seulement que ces
notes renvoient aux sections 9, 11, 12, 13, 18 et 90, et la note (24) à une
section dont le numéro surchargé se lit 69 ou 59, à propos des
« derivators ».}

\end{document}
